\makeatletter
\def\input@path{{./}{../}{../../}{../../../}{../../../../}{preamble/}{../preamble/}{../../preamble/}{../../../preamble/}}
\makeatother
\input{preamble_exams.tex}
\input{graphs/highlights.tex}
% \GTenHighlightsfalse % Uncomment to disable instructional highlights.
\input{graphs/specific_highlights.tex}
\COneTonePracticePalette

\setlength{\parskip}{3pt}\AtBeginDocument{\setlength{\abovedisplayskip}{4pt}\setlength{\belowdisplayskip}{4pt}}
\title{Grade 11 -- Term 1 \\ C9: Normal versus Student's t, with Errors}

\author{}\date{}
\begin{document}
\maketitle
\setcounter{tocdepth}{1}\setcounter{secnumdepth}{2}
\phantomsection\label{toc}\tableofcontents
\section{Models and calculation tools}
\label{sec:tools}
\Needspace{8\baselineskip}
\subsection*{Mathematical highlight guide}
\phantomsection\label{tone:C9:legend}
\PracticeHighlightLegend{G11-T1-C9}

% Relative to this PDF in 1-practice_activities; keep the sibling folder when hosting.
% Request a new window; the PDF viewer controls whether this becomes a tab.
\begin{tcolorbox}[colback=gray!2,colframe=gray!10,boxrule=0.7pt,
  arc=4pt,boxsep=3pt,left=8pt,right=8pt,top=5pt,bottom=5pt]
\centering
\href[pdfnewwindow=true]{https://nolopez338.github.io/global_economics/book/fundamentals/normal-table.pdf}{Open the normal and Student's t reference tables.}
\end{tcolorbox}
Page 1 gives $\ConeToneCDF{\Phi}(\ConeToneVariable{z})=\ConeToneProbability{P}(\ConeToneVariable{Z}\le \ConeToneVariable{z})$ for positive $\ConeToneVariable{z}$ on a 0.01 grid (through $2.59$).
Page 2 gives positive t critical values for upper-tail probabilities
$0.30,0.20,0.10,0.05,0.025,0.01,0.005$, for df $1$--$30$ and selected larger df.
Those entries do not give arbitrary CDF values. Use a suitable calculator's t CDF and inverse t CDF for this activity;
normal CDF and inverse normal are also allowed. Numerical CDF and quantile results in the solutions are supplementary to the tables.
Never turn table bounds into a single probability or error.

Write $\ConeToneCDF{F_\nu}(\ConeToneVariable{t})=\ConeToneProbability{P}(\ConeToneVariable{T}\le \ConeToneVariable{t})$ for $\ConeToneVariable{T}\sim \ConeToneVariable{t_\nu}$.
Symmetry gives $\ConeToneCDF{F_\nu}(-\ConeToneVariable{c})=\ConeToneCDF{1}-\ConeToneCDF{F_\nu}(\ConeToneVariable{c})$; a right tail is $1-\ConeToneCDF{F_\nu}(\ConeToneVariable{c})$.
An interval has probability $\ConeToneCDF{F_\nu}(\ConeToneVariable{b})-\ConeToneCDF{F_\nu}(\ConeToneVariable{\CeightToneScale{a}})$ and its outside complement has probability
$\ConeToneCDF{F_\nu}(\ConeToneVariable{\CeightToneScale{a}})+1-\ConeToneCDF{F_\nu}(\ConeToneVariable{b})$. Replace $\ConeToneCDF{F_\nu}$ by $\ConeToneCDF{\Phi}$ for normal calculations.
Endpoints have 0 probability. For central area $\ConeToneProbability{p}$, each tail has $(1-\ConeToneProbability{p})/2$;
for combined symmetric tails $\ConeToneVariable{q}$, each tail has $\ConeToneVariable{q}/2$.

A contextual location-scale model is explicitly $\ConeToneVariable{X}=\CfiveTonePopulationMean{m}+\CeightToneScale{a}\ConeToneVariable{T}$, with $\CeightToneScale{a}>0$ a \CeightToneScale{scale} \CeightToneScale{parameter};
$\CeightToneScale{a}$ is not automatically the standard deviation. Thus $\ConeToneVariable{T}=(\ConeToneVariable{X}-\CfiveTonePopulationMean{m})/\CeightToneScale{a}$. For a left cumulative probability $\ConeToneProbability{p}$, the contextual cutoff is $\ConeToneVariable{x}=\CfiveTonePopulationMean{m}+\CeightToneScale{a}\ConeToneCDF{F_\nu^{-1}}(\ConeToneProbability{p})$. A right-tail target $\ConeToneVariable{q}$ uses $\ConeToneProbability{p}=\ConeToneProbability{1}-\ConeToneVariable{q}$.
These are stipulated heavy-tailed models for signed deviations. Their parameters are given assumptions, not estimates from a sample. For $\CeightToneControl{\nu}>2$, their standard deviation is $\CeightToneScale{a}\sqrt{\CeightToneControl{\nu}/(\CeightToneControl{\nu}-2)}$, which differs from $\CeightToneScale{a}$.
For independent random observations from a normal population, with unknown \CfiveTonePopulationSD{population} \CfiveTonePopulationSD{standard} \CfiveTonePopulationSD{deviation},
the sampling statistic instead is
\[
\ConeToneVariable{T}=\frac{\CsevenToneSampleMean{\bar X}-\CfiveTonePopulationMean{\mu}}{\ConeToneVariable{S}/\sqrt \ConeToneUnknown{n}}\sim \ConeToneVariable{t_{n-1}}.
\]
Here $\CsevenToneSampleMean{\bar X}$ and $\ConeToneVariable{S}$ vary across samples; do not treat observed $\CsevenToneSampleSD{s}$ as a fixed \CfiveTonePopulationSD{population} \CfiveTonePopulationSD{standard} \CfiveTonePopulationSD{deviation}.

Attempt each task before checking the worked solution. Calculate with full precision, then report
probabilities to 4 decimals, standardized quantiles to 3 decimals, contextual values to 2 decimals,
and percentage errors to 2 decimals. Displayed rounded intermediate values are not the inputs used for errors.
Diagrams shade the requested region; curves continue beyond the displayed range.

\textbf{C9.} Compare normal approximations with t probabilities and quantiles, and interpret their errors.
For each comparison define $\ConeToneVariable{Z}\sim N(0,1)$ and $\ConeToneVariable{T}\sim \ConeToneVariable{t_\nu}$.
The normal calculation approximates the t calculation. In contextual units use
$\ConeToneVariable{X_N}=\CfiveTonePopulationMean{m}+\CeightToneScale{a}\ConeToneVariable{Z}$ and $\ConeToneVariable{X_T}=\CfiveTonePopulationMean{m}+\CeightToneScale{a}\ConeToneVariable{T}$, with exactly the same $\CfiveTonePopulationMean{m}$ and $\CeightToneScale{a}$.
Probability comparisons keep the same boundaries and event; inverse comparisons keep the same target probability. Normal cutoffs use $\ConeToneVariable{x_N}=\CfiveTonePopulationMean{m}+\CeightToneScale{a}\ConeToneCDF{\Phi^{-1}}(\ConeToneProbability{p})$; t cutoffs use $\ConeToneVariable{x_T}=\CfiveTonePopulationMean{m}+\CeightToneScale{a}\ConeToneCDF{F_\nu^{-1}}(\ConeToneProbability{p})$.
With normal result $\ConeToneVariable{r_N}$ and nonzero t reference $\ConeToneVariable{r_T}$, calculate
\[
 \CsixToneAbsoluteError{A}=|\ConeToneVariable{r_N}-\ConeToneVariable{r_T}|,\qquad \CsixToneRelativeError{R}=\frac{|\ConeToneVariable{r_N}-\ConeToneVariable{r_T}|}{|\ConeToneVariable{r_T}|},\qquad \CsixToneRelativeError{E}=\CsixToneRelativeError{100}\CsixToneRelativeError{R}\%.
\]
For inverse tasks, $\ConeToneVariable{r_T}$ and $\ConeToneVariable{r_N}$ are contextual cutoffs in the stated units. For intervals compare each endpoint separately. Probability absolute errors are measured on the 0--1 scale; $100\CsixToneAbsoluteError{A}$ is the difference in percentage points, whereas $100\CsixToneAbsoluteError{A}/|\ConeToneVariable{r_T}|$ is the \CsixToneRelativeError{percentage} \CsixToneRelativeError{error}. These errors depend on the quantity compared:
a standardized quantile and a contextual cutoff need not have the same \CsixToneRelativeError{relative} \CsixToneRelativeError{error}.
Heavier t tails give larger probabilities beyond positive cutoffs and more extreme quantiles at the same tail probability.
As df increases, t approaches normal. An interval entirely on 1 side of 0 need not contain less t probability.


\Needspace{8\baselineskip}
\subsection*{Notation reference}
\phantomsection\label{tone:C9:reference}
\COneToneDetailedReference{C9}
\section{Direct probability problems}\label{sec:c9-direct}
\SectionProblemsTableOfContents{
\SectionProblemLink{c9-d1}{Delivery schedule deviations} \\
\SectionProblemLink{c9-d2}{Production cost deviations} \\
\SectionProblemLink{c9-d3}{Electricity forecast deviations} \\
\SectionProblemLink{c9-d4}{Monthly sales deviations} \\
\SectionProblemLink{c9-d5}{Commute forecast deviations} \\
\SectionProblemLink{c9-d6}{Revenue planning alerts}}

\clearpage
\subsection{Delivery schedule deviations}
\label{c9-d1}
\textbf{Problem.} A courier uses a heavy-tailed planning model for delivery schedule deviations. Negative deviations mean early arrivals. Let $X$ be delivery time minus its scheduled time, measured in minutes. For this exercise, assume the location-scale model $X=0+5T$, where $T\sim t_{9}$ has 9 degrees of freedom. The location is $m=0$ and the scale is $a=5$ minutes; $a$ is a scale parameter, not automatically a standard deviation. Compare the reference model $X_T=0+5T$ with the normal approximation $X_N=0+5Z$, $Z\sim N(0,1)$, using identical location and scale.
\QuestionsTasks[sec:c9-direct]{{9}/{sol:c9-d1}/{Find the probability that a delivery arrives less than 7.5 minutes late (including early arrivals), and the probability that it arrives more than 6 minutes early. Standardize the boundaries, calculate the requested probabilities, and interpret each answer. Solve the same task under both models. Compare each probability or each contextual cutoff separately using absolute error and percentage error relative to t, calculated before rounding.}}
\subsubsection*{Worked solution}\phantomsection\label{sol:c9-d1}
\par\textbf{Translating the information into mathematics}\par
Let $\ConeToneVariable{X_T}$ and $\ConeToneVariable{X_N}$ denote delivery time minus its scheduled time, in minutes. The stated models specify $\ConeToneVariable{T}\sim \ConeToneVariable{t_{9}}$ and $\ConeToneVariable{Z}\sim N(0,1)$, and the location and scale specify $\ConeToneVariable{X_T}=\ConeToneVariable{0}+5\ConeToneVariable{T}$, $\ConeToneVariable{X_N}=\ConeToneVariable{0}+5\ConeToneVariable{Z}$. Thus $\CfiveTonePopulationMean{m}=\CfiveTonePopulationMean{0}$ and $\CeightToneScale{a}=\CeightToneScale{5}$; the scale is not a stipulated standard deviation. A contextual boundary $\ConeToneVariable{x}$ becomes $\ConeToneVariable{c}=(\ConeToneVariable{x}\ConeToneVariable{-0})/5$. The requested events are $\ConeToneProbability{P}(\ConeToneVariable{X}<7.5),\ \ConeToneProbability{P}(\ConeToneVariable{X}<-6)$, with $\ConeToneVariable{X}=\ConeToneVariable{X_T}$ or $\ConeToneVariable{X_N}$ for the appropriate model. Write the normal result as $\ConeToneProbability{\ConeToneVariable{r_N}}$ and the t reference as $\ConeToneProbability{\ConeToneVariable{r_T}}$. Compute $\CsixToneAbsoluteError{A}=|\ConeToneProbability{\ConeToneVariable{r_N}}-\ConeToneProbability{\ConeToneVariable{r_T}}|$, $\CsixToneRelativeError{R}=\CsixToneAbsoluteError{A}/|\ConeToneProbability{\ConeToneVariable{r_T}}|$, and $\CsixToneRelativeError{E}=\CsixToneRelativeError{100}\CsixToneRelativeError{R}\%$ with unrounded values.
\par\textbf{Mathematical reasoning and calculations}\par
Standardize each boundary using $(\ConeToneVariable{x}-0)/5$. In the task, $\ConeToneVariable{X}$ means $\ConeToneVariable{X_T}$ for the t model and $\ConeToneVariable{X_N}$ for the normal approximation.
\par\textbf{Part 1.}\[\ConeToneProbability{P}(\ConeToneVariable{X_T}<7.5)=\ConeToneProbability{P}(\ConeToneVariable{T}<\ConeToneProbability{1.5}).\]
\begin{center}
\begin{tikzpicture}[x=0.78cm,y=2.8cm]
\fill[COne Tone probability area] (-4.5,0) -- plot[domain=-4.50000000:1.50000000,samples=100] (\x,{0.3880349089*(1+\x*\x/9)^(-5)}) -- (1.5,0) -- cycle;
\draw[thick,COne Tone density curve] plot[domain=-4.50000000:4.50000000,samples=100] (\x,{0.3880349089*(1+\x*\x/9)^(-5)});
\draw[->] (-4.7,0)--(4.8,0);
\draw[densely dashed] (1.5,0)--(1.5,0.1271513);\node[below,font=\scriptsize] at (1.5,0) {$ 1.500 $};
\node[font=\small] at (0,0.49) {$ \ConeToneVariable{T}\sim \ConeToneVariable{t_{9}} $};
\node[below,font=\scriptsize] at (0,0) {$0$};\end{tikzpicture}\hfill
\begin{tikzpicture}[x=0.78cm,y=2.8cm]
\fill[COne Tone probability area] (-4.5,0) -- plot[domain=-4.50000000:1.50000000,samples=100] (\x,{0.3989422804*exp(-\x*\x/2)}) -- (1.5,0) -- cycle;
\draw[thick,COne Tone density curve] plot[domain=-4.50000000:4.50000000,samples=100] (\x,{0.3989422804*exp(-\x*\x/2)});
\draw[->] (-4.7,0)--(4.8,0);
\draw[densely dashed] (1.5,0)--(1.5,0.1295176);\node[below,font=\scriptsize] at (1.5,0) {$ 1.500 $};
\node[font=\small] at (0,0.49) {$ \ConeToneVariable{Z}\sim N(0,1) $};
\node[below,font=\scriptsize] at (0,0) {$0$};\end{tikzpicture}
\end{center}
For a negative boundary, use $\ConeToneCDF{F_\nu}(-\ConeToneVariable{c})=\ConeToneCDF{1}-\ConeToneCDF{F_\nu}(\ConeToneVariable{c})$ (and $\ConeToneCDF{\Phi}(-\ConeToneVariable{c})=\ConeToneCDF{1}-\ConeToneCDF{\Phi}(\ConeToneVariable{c})$). The region is cumulative area to the left.
Supplementary CDF values: $\ConeToneCDF{F_{9}}(1.5)\approx \ConeToneCDF{0.9160747}$.
\[\ConeToneProbability{P}(\ConeToneVariable{T}<1.5)=\ConeToneCDF{F_{9}}(\ConeToneProbability{1.5})\approx \ConeToneProbability{0.9161}.\]
\[\ConeToneProbability{P}(\ConeToneVariable{Z}<1.5)=\ConeToneCDF{\Phi}(\ConeToneProbability{1.5})\approx \ConeToneProbability{0.9332}.\]
\[\CsixToneAbsoluteError{A}=|\CsixToneAbsoluteError{0.9331928}-(0.9160747)|\approx \CsixToneAbsoluteError{0.017118},\quad \CsixToneRelativeError{R}=\frac{\CsixToneAbsoluteError{A}}{\CeightToneScale{0.9160747}}\approx \CsixToneRelativeError{0.018686},\quad \CsixToneRelativeError{E}\approx \CsixToneRelativeError{1.87\%}.\]
\textbf{Interpretation.} Under t, the probability that a delivery arrives less than 7.5 minutes late is 91.61\%. The normal approximation overestimates that probability, by 1.71 percentage points (1.87\% relative to t).
\par\textbf{Part 2.}\[\ConeToneProbability{P}(\ConeToneVariable{X_T}<-6)=\ConeToneProbability{P}(\ConeToneVariable{T}<\ConeToneProbability{-1.2}).\]
\begin{center}
\begin{tikzpicture}[x=0.78cm,y=2.8cm]
\fill[COne Tone probability area] (-4.5,0) -- plot[domain=-4.50000000:-1.20000000,samples=100] (\x,{0.3880349089*(1+\x*\x/9)^(-5)}) -- (-1.2,0) -- cycle;
\draw[thick,COne Tone density curve] plot[domain=-4.50000000:4.50000000,samples=100] (\x,{0.3880349089*(1+\x*\x/9)^(-5)});
\draw[->] (-4.7,0)--(4.8,0);
\draw[densely dashed] (-1.2,0)--(-1.2,0.1847485);\node[below,font=\scriptsize] at (-1.2,0) {$ -1.200 $};
\node[font=\small] at (0,0.49) {$ \ConeToneVariable{T}\sim \ConeToneVariable{t_{9}} $};
\node[below,font=\scriptsize] at (0,0) {$0$};\end{tikzpicture}\hfill
\begin{tikzpicture}[x=0.78cm,y=2.8cm]
\fill[COne Tone probability area] (-4.5,0) -- plot[domain=-4.50000000:-1.20000000,samples=100] (\x,{0.3989422804*exp(-\x*\x/2)}) -- (-1.2,0) -- cycle;
\draw[thick,COne Tone density curve] plot[domain=-4.50000000:4.50000000,samples=100] (\x,{0.3989422804*exp(-\x*\x/2)});
\draw[->] (-4.7,0)--(4.8,0);
\draw[densely dashed] (-1.2,0)--(-1.2,0.1941861);\node[below,font=\scriptsize] at (-1.2,0) {$ -1.200 $};
\node[font=\small] at (0,0.49) {$ \ConeToneVariable{Z}\sim N(0,1) $};
\node[below,font=\scriptsize] at (0,0) {$0$};\end{tikzpicture}
\end{center}
For a negative boundary, use $\ConeToneCDF{F_\nu}(-\ConeToneVariable{c})=\ConeToneCDF{1}-\ConeToneCDF{F_\nu}(\ConeToneVariable{c})$ (and $\ConeToneCDF{\Phi}(-\ConeToneVariable{c})=\ConeToneCDF{1}-\ConeToneCDF{\Phi}(\ConeToneVariable{c})$). The region is cumulative area to the left.
By symmetry, $\ConeToneCDF{F_{9}}(-1.2)=\ConeToneCDF{1}-\ConeToneCDF{F_{9}}(1.2)\approx \ConeToneCDF{1}-0.8696134=\ConeToneCDF{0.1303866}$. Supplementary CDF values: $\ConeToneCDF{F_{9}}(-1.2)\approx \ConeToneCDF{0.1303866}$.
\[\ConeToneProbability{P}(\ConeToneVariable{T}<-1.2)=\ConeToneCDF{F_{9}}(\ConeToneProbability{-1.2})\approx \ConeToneProbability{0.1304}.\]
\[\ConeToneProbability{P}(\ConeToneVariable{Z}<-1.2)=\ConeToneCDF{\Phi}(\ConeToneProbability{-1.2})\approx \ConeToneProbability{0.1151}.\]
\[\CsixToneAbsoluteError{A}=|\CsixToneAbsoluteError{0.1150697}-(0.1303866)|\approx \CsixToneAbsoluteError{0.015317},\quad \CsixToneRelativeError{R}=\frac{\CsixToneAbsoluteError{A}}{\CeightToneScale{0.1303866}}\approx \CsixToneRelativeError{0.117473},\quad \CsixToneRelativeError{E}\approx \CsixToneRelativeError{11.75\%}.\]
\textbf{Interpretation.} Under t, the probability that a delivery arrives more than 6 minutes early is 13.04\%. The normal approximation underestimates that probability, by 1.53 percentage points (11.75\% relative to t).

\secInicio[c9-d1][sec:c9-direct]

\clearpage
\subsection{Production cost deviations}
\label{c9-d2}
\textbf{Problem.} A manufacturer models batch cost deviations for budget planning. Negative deviations represent savings. Let $X$ be batch cost minus its budget, measured in monetary units. For this exercise, assume the location-scale model $X=0+20T$, where $T\sim t_{14}$ has 14 degrees of freedom. The location is $m=0$ and the scale is $a=20$ monetary units; $a$ is a scale parameter, not automatically a standard deviation. Compare the reference model $X_T=0+20T$ with the normal approximation $X_N=0+20Z$, $Z\sim N(0,1)$, using identical location and scale.
\QuestionsTasks[sec:c9-direct]{{9}/{sol:c9-d2}/{Find the probability that batch cost exceeds budget by more than 30 monetary units, and the probability that it exceeds the budget reduced by 16 monetary units. Standardize the boundaries, calculate the requested probabilities, and interpret each answer. Solve the same task under both models. Compare each probability or each contextual cutoff separately using absolute error and percentage error relative to t, calculated before rounding.}}
\subsubsection*{Worked solution}\phantomsection\label{sol:c9-d2}
\par\textbf{Translating the information into mathematics}\par
Let $\ConeToneVariable{X_T}$ and $\ConeToneVariable{X_N}$ denote batch cost minus its budget, in monetary units. The stated models specify $\ConeToneVariable{T}\sim \ConeToneVariable{t_{14}}$ and $\ConeToneVariable{Z}\sim N(0,1)$, and the location and scale specify $\ConeToneVariable{X_T}=\ConeToneVariable{0}+20\ConeToneVariable{T}$, $\ConeToneVariable{X_N}=\ConeToneVariable{0}+20\ConeToneVariable{Z}$. Thus $\CfiveTonePopulationMean{m}=\CfiveTonePopulationMean{0}$ and $\CeightToneScale{a}=\CeightToneScale{20}$; the scale is not a stipulated standard deviation. A contextual boundary $\ConeToneVariable{x}$ becomes $\ConeToneVariable{c}=(\ConeToneVariable{x}\ConeToneVariable{-0})/20$. The requested events are $\ConeToneProbability{P}(\ConeToneVariable{X}>30),\ \ConeToneProbability{P}(\ConeToneVariable{X}>-16)$, with $\ConeToneVariable{X}=\ConeToneVariable{X_T}$ or $\ConeToneVariable{X_N}$ for the appropriate model. Write the normal result as $\ConeToneProbability{\ConeToneVariable{r_N}}$ and the t reference as $\ConeToneProbability{\ConeToneVariable{r_T}}$. Compute $\CsixToneAbsoluteError{A}=|\ConeToneProbability{\ConeToneVariable{r_N}}-\ConeToneProbability{\ConeToneVariable{r_T}}|$, $\CsixToneRelativeError{R}=\CsixToneAbsoluteError{A}/|\ConeToneProbability{\ConeToneVariable{r_T}}|$, and $\CsixToneRelativeError{E}=\CsixToneRelativeError{100}\CsixToneRelativeError{R}\%$ with unrounded values.
\par\textbf{Mathematical reasoning and calculations}\par
Standardize each boundary using $(\ConeToneVariable{x}-0)/20$. In the task, $\ConeToneVariable{X}$ means $\ConeToneVariable{X_T}$ for the t model and $\ConeToneVariable{X_N}$ for the normal approximation.
\par\textbf{Part 1.}\[\ConeToneProbability{P}(\ConeToneVariable{X_T}>30)=\ConeToneProbability{P}(\ConeToneVariable{T}>\ConeToneProbability{1.5}).\]
\begin{center}
\begin{tikzpicture}[x=0.78cm,y=2.8cm]
\fill[COne Tone probability area] (1.5,0) -- plot[domain=1.50000000:4.50000000,samples=100] (\x,{0.3918874558*(1+\x*\x/14)^(-7.5)}) -- (4.5,0) -- cycle;
\draw[thick,COne Tone density curve] plot[domain=-4.50000000:4.50000000,samples=100] (\x,{0.3918874558*(1+\x*\x/14)^(-7.5)});
\draw[->] (-4.7,0)--(4.8,0);
\draw[densely dashed] (1.5,0)--(1.5,0.1281508);\node[below,font=\scriptsize] at (1.5,0) {$ 1.500 $};
\node[font=\small] at (0,0.49) {$ \ConeToneVariable{T}\sim \ConeToneVariable{t_{14}} $};
\node[below,font=\scriptsize] at (0,0) {$0$};\end{tikzpicture}\hfill
\begin{tikzpicture}[x=0.78cm,y=2.8cm]
\fill[COne Tone probability area] (1.5,0) -- plot[domain=1.50000000:4.50000000,samples=100] (\x,{0.3989422804*exp(-\x*\x/2)}) -- (4.5,0) -- cycle;
\draw[thick,COne Tone density curve] plot[domain=-4.50000000:4.50000000,samples=100] (\x,{0.3989422804*exp(-\x*\x/2)});
\draw[->] (-4.7,0)--(4.8,0);
\draw[densely dashed] (1.5,0)--(1.5,0.1295176);\node[below,font=\scriptsize] at (1.5,0) {$ 1.500 $};
\node[font=\small] at (0,0.49) {$ \ConeToneVariable{Z}\sim N(0,1) $};
\node[below,font=\scriptsize] at (0,0) {$0$};\end{tikzpicture}
\end{center}
For a negative boundary, use $\ConeToneCDF{F_\nu}(-\ConeToneVariable{c})=\ConeToneCDF{1}-\ConeToneCDF{F_\nu}(\ConeToneVariable{c})$ (and $\ConeToneCDF{\Phi}(-\ConeToneVariable{c})=\ConeToneCDF{1}-\ConeToneCDF{\Phi}(\ConeToneVariable{c})$). Take the complement of the cumulative area.
Supplementary CDF values: $\ConeToneCDF{F_{14}}(1.5)\approx \ConeToneCDF{0.9220873}$.
\[\ConeToneProbability{P}(\ConeToneVariable{T}>1.5)=\ConeToneProbability{1}-\ConeToneCDF{F_{14}}(1.5)\approx \ConeToneProbability{0.0779}.\]
\[\ConeToneProbability{P}(\ConeToneVariable{Z}>1.5)=\ConeToneProbability{1}-\ConeToneCDF{\Phi}(1.5)\approx \ConeToneProbability{0.0668}.\]
\[\CsixToneAbsoluteError{A}=|\CsixToneAbsoluteError{0.0668072}-(0.0779127)|\approx \CsixToneAbsoluteError{0.011105},\quad \CsixToneRelativeError{R}=\frac{\CsixToneAbsoluteError{A}}{\CeightToneScale{0.0779127}}\approx \CsixToneRelativeError{0.142537},\quad \CsixToneRelativeError{E}\approx \CsixToneRelativeError{14.25\%}.\]
\textbf{Interpretation.} Under t, the probability that a batch costs more than 30 monetary units above budget is 7.79\%. The normal approximation underestimates that probability, by 1.11 percentage points (14.25\% relative to t).
\par\textbf{Part 2.}\[\ConeToneProbability{P}(\ConeToneVariable{X_T}>-16)=\ConeToneProbability{P}(\ConeToneVariable{T}>\ConeToneProbability{-0.8}).\]
\begin{center}
\begin{tikzpicture}[x=0.78cm,y=2.8cm]
\fill[COne Tone probability area] (-0.8,0) -- plot[domain=-0.80000000:4.50000000,samples=100] (\x,{0.3918874558*(1+\x*\x/14)^(-7.5)}) -- (4.5,0) -- cycle;
\draw[thick,COne Tone density curve] plot[domain=-4.50000000:4.50000000,samples=100] (\x,{0.3918874558*(1+\x*\x/14)^(-7.5)});
\draw[->] (-4.7,0)--(4.8,0);
\draw[densely dashed] (-0.8,0)--(-0.8,0.2802616);\node[below,font=\scriptsize] at (-0.8,0) {$ -0.800 $};
\node[font=\small] at (0,0.49) {$ \ConeToneVariable{T}\sim \ConeToneVariable{t_{14}} $};
\node[below,font=\scriptsize] at (0,0) {$0$};\end{tikzpicture}\hfill
\begin{tikzpicture}[x=0.78cm,y=2.8cm]
\fill[COne Tone probability area] (-0.8,0) -- plot[domain=-0.80000000:4.50000000,samples=100] (\x,{0.3989422804*exp(-\x*\x/2)}) -- (4.5,0) -- cycle;
\draw[thick,COne Tone density curve] plot[domain=-4.50000000:4.50000000,samples=100] (\x,{0.3989422804*exp(-\x*\x/2)});
\draw[->] (-4.7,0)--(4.8,0);
\draw[densely dashed] (-0.8,0)--(-0.8,0.2896916);\node[below,font=\scriptsize] at (-0.8,0) {$ -0.800 $};
\node[font=\small] at (0,0.49) {$ \ConeToneVariable{Z}\sim N(0,1) $};
\node[below,font=\scriptsize] at (0,0) {$0$};\end{tikzpicture}
\end{center}
For a negative boundary, use $\ConeToneCDF{F_\nu}(-\ConeToneVariable{c})=\ConeToneCDF{1}-\ConeToneCDF{F_\nu}(\ConeToneVariable{c})$ (and $\ConeToneCDF{\Phi}(-\ConeToneVariable{c})=\ConeToneCDF{1}-\ConeToneCDF{\Phi}(\ConeToneVariable{c})$). Take the complement of the cumulative area.
By symmetry, $\ConeToneCDF{F_{14}}(-0.8)=\ConeToneCDF{1}-\ConeToneCDF{F_{14}}(0.8)\approx \ConeToneCDF{1}-0.7814622=\ConeToneCDF{0.2185378}$. Supplementary CDF values: $\ConeToneCDF{F_{14}}(-0.8)\approx \ConeToneCDF{0.2185378}$.
\[\ConeToneProbability{P}(\ConeToneVariable{T}>-0.8)=\ConeToneProbability{1}-\ConeToneCDF{F_{14}}(-0.8)\approx \ConeToneProbability{0.7815}.\]
\[\ConeToneProbability{P}(\ConeToneVariable{Z}>-0.8)=\ConeToneProbability{1}-\ConeToneCDF{\Phi}(-0.8)\approx \ConeToneProbability{0.7881}.\]
\[\CsixToneAbsoluteError{A}=|\CsixToneAbsoluteError{0.7881446}-(0.7814622)|\approx \CsixToneAbsoluteError{0.006682},\quad \CsixToneRelativeError{R}=\frac{\CsixToneAbsoluteError{A}}{\CeightToneScale{0.7814622}}\approx \CsixToneRelativeError{0.008551},\quad \CsixToneRelativeError{E}\approx \CsixToneRelativeError{0.86\%}.\]
\textbf{Interpretation.} Under t, the probability that a batch costs more than the budget reduced by 16 monetary units is 78.15\%. The normal approximation overestimates that probability, by 0.67 percentage points (0.86\% relative to t).

\secInicio[c9-d2][sec:c9-direct]

\clearpage
\subsection{Electricity forecast deviations}
\label{c9-d3}
\textbf{Problem.} An energy planner models household electricity forecast errors. Negative errors mean use below forecast. Let $X$ be household electricity use minus its forecast, measured in kilowatt-hours. For this exercise, assume the location-scale model $X=0+50T$, where $T\sim t_{19}$ has 19 degrees of freedom. The location is $m=0$ and the scale is $a=50$ kilowatt-hours; $a$ is a scale parameter, not automatically a standard deviation. Compare the reference model $X_T=0+50T$ with the normal approximation $X_N=0+50Z$, $Z\sim N(0,1)$, using identical location and scale.
\QuestionsTasks[sec:c9-direct]{{9}/{sol:c9-d3}/{Find the probability that household electricity use is between 12.5 and 57.5 kilowatt-hours above forecast. Standardize the boundaries, calculate the requested probabilities, and interpret each answer. Solve the same task under both models. Compare each probability or each contextual cutoff separately using absolute error and percentage error relative to t, calculated before rounding.}}
\subsubsection*{Worked solution}\phantomsection\label{sol:c9-d3}
\par\textbf{Translating the information into mathematics}\par
Let $\ConeToneVariable{X_T}$ and $\ConeToneVariable{X_N}$ denote household electricity use minus its forecast, in kilowatt-hours. The stated models specify $\ConeToneVariable{T}\sim \ConeToneVariable{t_{19}}$ and $\ConeToneVariable{Z}\sim N(0,1)$, and the location and scale specify $\ConeToneVariable{X_T}=\ConeToneVariable{0}+50\ConeToneVariable{T}$, $\ConeToneVariable{X_N}=\ConeToneVariable{0}+50\ConeToneVariable{Z}$. Thus $\CfiveTonePopulationMean{m}=\CfiveTonePopulationMean{0}$ and $\CeightToneScale{a}=\CeightToneScale{50}$; the scale is not a stipulated standard deviation. A contextual boundary $\ConeToneVariable{x}$ becomes $\ConeToneVariable{c}=(\ConeToneVariable{x}\ConeToneVariable{-0})/50$. The requested events are $\ConeToneProbability{P}(12.5<\ConeToneVariable{X}<57.5)$, with $\ConeToneVariable{X}=\ConeToneVariable{X_T}$ or $\ConeToneVariable{X_N}$ for the appropriate model. Write the normal result as $\ConeToneProbability{\ConeToneVariable{r_N}}$ and the t reference as $\ConeToneProbability{\ConeToneVariable{r_T}}$. Compute $\CsixToneAbsoluteError{A}=|\ConeToneProbability{\ConeToneVariable{r_N}}-\ConeToneProbability{\ConeToneVariable{r_T}}|$, $\CsixToneRelativeError{R}=\CsixToneAbsoluteError{A}/|\ConeToneProbability{\ConeToneVariable{r_T}}|$, and $\CsixToneRelativeError{E}=\CsixToneRelativeError{100}\CsixToneRelativeError{R}\%$ with unrounded values.
\par\textbf{Mathematical reasoning and calculations}\par
Standardize each boundary using $(\ConeToneVariable{x}-0)/50$. In the task, $\ConeToneVariable{X}$ means $\ConeToneVariable{X_T}$ for the t model and $\ConeToneVariable{X_N}$ for the normal approximation.
\par\textbf{Part 1.}\[\ConeToneProbability{P}(12.5<\ConeToneVariable{X_T}<57.5)=\ConeToneProbability{P}(\ConeToneProbability{0.25}<\ConeToneVariable{T}<1.15).\]
\begin{center}
\begin{tikzpicture}[x=0.78cm,y=2.8cm]
\fill[COne Tone probability area] (0.25,0) -- plot[domain=0.25000000:1.15000000,samples=100] (\x,{0.3937298073*(1+\x*\x/19)^(-10)}) -- (1.15,0) -- cycle;
\draw[thick,COne Tone density curve] plot[domain=-4.50000000:4.50000000,samples=100] (\x,{0.3937298073*(1+\x*\x/19)^(-10)});
\draw[->] (-4.7,0)--(4.8,0);
\draw[densely dashed] (0.25,0)--(0.25,0.3810094);\node[below,font=\scriptsize] at (0.25,0) {$ 0.250 $};
\draw[densely dashed] (1.15,0)--(1.15,0.2008922);\node[below,font=\scriptsize] at (1.15,0) {$ 1.150 $};
\node[font=\small] at (0,0.49) {$ \ConeToneVariable{T}\sim \ConeToneVariable{t_{19}} $};
\node[below,font=\scriptsize] at (0,0) {$0$};\end{tikzpicture}\hfill
\begin{tikzpicture}[x=0.78cm,y=2.8cm]
\fill[COne Tone probability area] (0.25,0) -- plot[domain=0.25000000:1.15000000,samples=100] (\x,{0.3989422804*exp(-\x*\x/2)}) -- (1.15,0) -- cycle;
\draw[thick,COne Tone density curve] plot[domain=-4.50000000:4.50000000,samples=100] (\x,{0.3989422804*exp(-\x*\x/2)});
\draw[->] (-4.7,0)--(4.8,0);
\draw[densely dashed] (0.25,0)--(0.25,0.3866681);\node[below,font=\scriptsize] at (0.25,0) {$ 0.250 $};
\draw[densely dashed] (1.15,0)--(1.15,0.2059363);\node[below,font=\scriptsize] at (1.15,0) {$ 1.150 $};
\node[font=\small] at (0,0.49) {$ \ConeToneVariable{Z}\sim N(0,1) $};
\node[below,font=\scriptsize] at (0,0) {$0$};\end{tikzpicture}
\end{center}
For a negative boundary, use $\ConeToneCDF{F_\nu}(-\ConeToneVariable{c})=\ConeToneCDF{1}-\ConeToneCDF{F_\nu}(\ConeToneVariable{c})$ (and $\ConeToneCDF{\Phi}(-\ConeToneVariable{c})=\ConeToneCDF{1}-\ConeToneCDF{\Phi}(\ConeToneVariable{c})$). Subtract the left cumulative area from the right cumulative area.
Supplementary CDF values: $\ConeToneCDF{F_{19}}(0.25)\approx \ConeToneCDF{0.5973648},\quad \ConeToneCDF{F_{19}}(1.15)\approx \ConeToneCDF{0.8677933}$.
\[\ConeToneProbability{P}(0.25<\ConeToneVariable{T}<1.15)=\ConeToneCDF{F_{19}}(\ConeToneProbability{1.15})-\ConeToneCDF{F_{19}}(0.25)\approx \ConeToneProbability{0.2704}.\]
\[\ConeToneProbability{P}(0.25<\ConeToneVariable{Z}<1.15)=\ConeToneCDF{\Phi}(\ConeToneProbability{1.15})-\ConeToneCDF{\Phi}(0.25)\approx \ConeToneProbability{0.2762}.\]
\[\CsixToneAbsoluteError{A}=|\CsixToneAbsoluteError{0.2762217}-(0.2704285)|\approx \CsixToneAbsoluteError{0.005793},\quad \CsixToneRelativeError{R}=\frac{\CsixToneAbsoluteError{A}}{\CeightToneScale{0.2704285}}\approx \CsixToneRelativeError{0.021422},\quad \CsixToneRelativeError{E}\approx \CsixToneRelativeError{2.14\%}.\]
\textbf{Interpretation.} Under t, the probability that household use is between 12.5 and 57.5 kWh above forecast is 27.04\%. The normal approximation overestimates that probability, by 0.58 percentage points (2.14\% relative to t).

\secInicio[c9-d3][sec:c9-direct]

\clearpage
\subsection{Monthly sales deviations}
\label{c9-d4}
\textbf{Problem.} A retailer uses a continuous model for monthly sales deviations when planning inventory. Negative deviations mean sales below target. Let $X$ be monthly sales minus its target, measured in units. For this exercise, assume the location-scale model $X=0+20T$, where $T\sim t_{11}$ has 11 degrees of freedom. The location is $m=0$ and the scale is $a=20$ units; $a$ is a scale parameter, not automatically a standard deviation. Compare the reference model $X_T=0+20T$ with the normal approximation $X_N=0+20Z$, $Z\sim N(0,1)$, using identical location and scale.
\QuestionsTasks[sec:c9-direct]{{9}/{sol:c9-d4}/{Find the probability that monthly sales fall short of target by between 17 and 36 units. Standardize the boundaries, calculate the requested probabilities, and interpret each answer. Solve the same task under both models. Compare each probability or each contextual cutoff separately using absolute error and percentage error relative to t, calculated before rounding.}}
\subsubsection*{Worked solution}\phantomsection\label{sol:c9-d4}
\par\textbf{Translating the information into mathematics}\par
Let $\ConeToneVariable{X_T}$ and $\ConeToneVariable{X_N}$ denote monthly sales minus its target, in units sold. The stated models specify $\ConeToneVariable{T}\sim \ConeToneVariable{t_{11}}$ and $\ConeToneVariable{Z}\sim N(0,1)$, and the location and scale specify $\ConeToneVariable{X_T}=\ConeToneVariable{0}+20\ConeToneVariable{T}$, $\ConeToneVariable{X_N}=\ConeToneVariable{0}+20\ConeToneVariable{Z}$. Thus $\CfiveTonePopulationMean{m}=\CfiveTonePopulationMean{0}$ and $\CeightToneScale{a}=\CeightToneScale{20}$; the scale is not a stipulated standard deviation. A contextual boundary $\ConeToneVariable{x}$ becomes $\ConeToneVariable{c}=(\ConeToneVariable{x}\ConeToneVariable{-0})/20$. The requested events are $\ConeToneProbability{P}(-36<\ConeToneVariable{X}<-17)$, with $\ConeToneVariable{X}=\ConeToneVariable{X_T}$ or $\ConeToneVariable{X_N}$ for the appropriate model. Write the normal result as $\ConeToneProbability{\ConeToneVariable{r_N}}$ and the t reference as $\ConeToneProbability{\ConeToneVariable{r_T}}$. Compute $\CsixToneAbsoluteError{A}=|\ConeToneProbability{\ConeToneVariable{r_N}}-\ConeToneProbability{\ConeToneVariable{r_T}}|$, $\CsixToneRelativeError{R}=\CsixToneAbsoluteError{A}/|\ConeToneProbability{\ConeToneVariable{r_T}}|$, and $\CsixToneRelativeError{E}=\CsixToneRelativeError{100}\CsixToneRelativeError{R}\%$ with unrounded values.
\par\textbf{Mathematical reasoning and calculations}\par
Standardize each boundary using $(\ConeToneVariable{x}-0)/20$. In the task, $\ConeToneVariable{X}$ means $\ConeToneVariable{X_T}$ for the t model and $\ConeToneVariable{X_N}$ for the normal approximation.
\par\textbf{Part 1.}\[\ConeToneProbability{P}(-36<\ConeToneVariable{X_T}<-17)=\ConeToneProbability{P}(\ConeToneProbability{-1.8}<\ConeToneVariable{T}<-0.85).\]
\begin{center}
\begin{tikzpicture}[x=0.78cm,y=2.8cm]
\fill[COne Tone probability area] (-1.8,0) -- plot[domain=-1.80000000:-0.85000000,samples=100] (\x,{0.3899897571*(1+\x*\x/11)^(-6)}) -- (-0.85,0) -- cycle;
\draw[thick,COne Tone density curve] plot[domain=-4.50000000:4.50000000,samples=100] (\x,{0.3899897571*(1+\x*\x/11)^(-6)});
\draw[->] (-4.7,0)--(4.8,0);
\draw[densely dashed] (-1.8,0)--(-1.8,0.0828608);\node[below,font=\scriptsize] at (-1.8,0) {$ -1.800 $};
\draw[densely dashed] (-0.85,0)--(-0.85,0.2662489);\node[below,font=\scriptsize] at (-0.85,0) {$ -0.850 $};
\node[font=\small] at (0,0.49) {$ \ConeToneVariable{T}\sim \ConeToneVariable{t_{11}} $};
\node[below,font=\scriptsize] at (0,0) {$0$};\end{tikzpicture}\hfill
\begin{tikzpicture}[x=0.78cm,y=2.8cm]
\fill[COne Tone probability area] (-1.8,0) -- plot[domain=-1.80000000:-0.85000000,samples=100] (\x,{0.3989422804*exp(-\x*\x/2)}) -- (-0.85,0) -- cycle;
\draw[thick,COne Tone density curve] plot[domain=-4.50000000:4.50000000,samples=100] (\x,{0.3989422804*exp(-\x*\x/2)});
\draw[->] (-4.7,0)--(4.8,0);
\draw[densely dashed] (-1.8,0)--(-1.8,0.0789502);\node[below,font=\scriptsize] at (-1.8,0) {$ -1.800 $};
\draw[densely dashed] (-0.85,0)--(-0.85,0.2779849);\node[below,font=\scriptsize] at (-0.85,0) {$ -0.850 $};
\node[font=\small] at (0,0.49) {$ \ConeToneVariable{Z}\sim N(0,1) $};
\node[below,font=\scriptsize] at (0,0) {$0$};\end{tikzpicture}
\end{center}
For a negative boundary, use $\ConeToneCDF{F_\nu}(-\ConeToneVariable{c})=\ConeToneCDF{1}-\ConeToneCDF{F_\nu}(\ConeToneVariable{c})$ (and $\ConeToneCDF{\Phi}(-\ConeToneVariable{c})=\ConeToneCDF{1}-\ConeToneCDF{\Phi}(\ConeToneVariable{c})$). Subtract the left cumulative area from the right cumulative area.
By symmetry, $\ConeToneCDF{F_{11}}(-1.8)=\ConeToneCDF{1}-\ConeToneCDF{F_{11}}(1.8)\approx \ConeToneCDF{1}-0.9503421=\ConeToneCDF{0.0496579}$. By symmetry, $\ConeToneCDF{F_{11}}(-0.85)=\ConeToneCDF{1}-\ConeToneCDF{F_{11}}(0.85)\approx \ConeToneCDF{1}-0.7932783=\ConeToneCDF{0.2067217}$. Supplementary CDF values: $\ConeToneCDF{F_{11}}(-1.8)\approx \ConeToneCDF{0.0496579},\quad \ConeToneCDF{F_{11}}(-0.85)\approx \ConeToneCDF{0.2067217}$.
\[\ConeToneProbability{P}(-1.8<\ConeToneVariable{T}<-0.85)=\ConeToneCDF{F_{11}}(\ConeToneProbability{-0.85})-\ConeToneCDF{F_{11}}(-1.8)\approx \ConeToneProbability{0.1571}.\]
\[\ConeToneProbability{P}(-1.8<\ConeToneVariable{Z}<-0.85)=\ConeToneCDF{\Phi}(\ConeToneProbability{-0.85})-\ConeToneCDF{\Phi}(-1.8)\approx \ConeToneProbability{0.1617}.\]
\[\CsixToneAbsoluteError{A}=|\CsixToneAbsoluteError{0.1617322}-(0.1570638)|\approx \CsixToneAbsoluteError{0.004668},\quad \CsixToneRelativeError{R}=\frac{\CsixToneAbsoluteError{A}}{\CeightToneScale{0.1570638}}\approx \CsixToneRelativeError{0.029723},\quad \CsixToneRelativeError{E}\approx \CsixToneRelativeError{2.97\%}.\]
\textbf{Interpretation.} Under t, the probability that monthly sales fall between 17 and 36 units below target is 15.71\%. The normal approximation overestimates that probability, by 0.47 percentage points (2.97\% relative to t).

\secInicio[c9-d4][sec:c9-direct]

\clearpage
\subsection{Commute forecast deviations}
\label{c9-d5}
\textbf{Problem.} A transport planner models errors in commute-time forecasts. Negative errors mean a shorter trip than forecast. Let $X$ be commute time minus its forecast, measured in minutes. For this exercise, assume the location-scale model $X=0+10T$, where $T\sim t_{24}$ has 24 degrees of freedom. The location is $m=0$ and the scale is $a=10$ minutes; $a$ is a scale parameter, not automatically a standard deviation. Compare the reference model $X_T=0+10T$ with the normal approximation $X_N=0+10Z$, $Z\sim N(0,1)$, using identical location and scale.
\QuestionsTasks[sec:c9-direct]{{9}/{sol:c9-d5}/{Find the probability that commute time is between 9.5 minutes below forecast and 10.5 minutes above forecast. Standardize the boundaries, calculate the requested probabilities, and interpret each answer. Solve the same task under both models. Compare each probability or each contextual cutoff separately using absolute error and percentage error relative to t, calculated before rounding.}}
\subsubsection*{Worked solution}\phantomsection\label{sol:c9-d5}
\par\textbf{Translating the information into mathematics}\par
Let $\ConeToneVariable{X_T}$ and $\ConeToneVariable{X_N}$ denote commute time minus its forecast, in minutes. The stated models specify $\ConeToneVariable{T}\sim \ConeToneVariable{t_{24}}$ and $\ConeToneVariable{Z}\sim N(0,1)$, and the location and scale specify $\ConeToneVariable{X_T}=\ConeToneVariable{0}+10\ConeToneVariable{T}$, $\ConeToneVariable{X_N}=\ConeToneVariable{0}+10\ConeToneVariable{Z}$. Thus $\CfiveTonePopulationMean{m}=\CfiveTonePopulationMean{0}$ and $\CeightToneScale{a}=\CeightToneScale{10}$; the scale is not a stipulated standard deviation. A contextual boundary $\ConeToneVariable{x}$ becomes $\ConeToneVariable{c}=(\ConeToneVariable{x}\ConeToneVariable{-0})/10$. The requested events are $\ConeToneProbability{P}(-9.5<\ConeToneVariable{X}<10.5)$, with $\ConeToneVariable{X}=\ConeToneVariable{X_T}$ or $\ConeToneVariable{X_N}$ for the appropriate model. Write the normal result as $\ConeToneProbability{\ConeToneVariable{r_N}}$ and the t reference as $\ConeToneProbability{\ConeToneVariable{r_T}}$. Compute $\CsixToneAbsoluteError{A}=|\ConeToneProbability{\ConeToneVariable{r_N}}-\ConeToneProbability{\ConeToneVariable{r_T}}|$, $\CsixToneRelativeError{R}=\CsixToneAbsoluteError{A}/|\ConeToneProbability{\ConeToneVariable{r_T}}|$, and $\CsixToneRelativeError{E}=\CsixToneRelativeError{100}\CsixToneRelativeError{R}\%$ with unrounded values.
\par\textbf{Mathematical reasoning and calculations}\par
Standardize each boundary using $(\ConeToneVariable{x}-0)/10$. In the task, $\ConeToneVariable{X}$ means $\ConeToneVariable{X_T}$ for the t model and $\ConeToneVariable{X_N}$ for the normal approximation.
\par\textbf{Part 1.}\[\ConeToneProbability{P}(-9.5<\ConeToneVariable{X_T}<10.5)=\ConeToneProbability{P}(\ConeToneProbability{-0.95}<\ConeToneVariable{T}<1.05).\]
\begin{center}
\begin{tikzpicture}[x=0.78cm,y=2.8cm]
\fill[COne Tone probability area] (-0.95,0) -- plot[domain=-0.95000000:1.05000000,samples=100] (\x,{0.3948093882*(1+\x*\x/24)^(-12.5)}) -- (1.05,0) -- cycle;
\draw[thick,COne Tone density curve] plot[domain=-4.50000000:4.50000000,samples=100] (\x,{0.3948093882*(1+\x*\x/24)^(-12.5)});
\draw[->] (-4.7,0)--(4.8,0);
\draw[densely dashed] (-0.95,0)--(-0.95,0.2488807);\node[below,font=\scriptsize] at (-0.95,0) {$ -0.950 $};
\draw[densely dashed] (1.05,0)--(1.05,0.2251987);\node[below,font=\scriptsize] at (1.05,0) {$ 1.050 $};
\node[font=\small] at (0,0.49) {$ \ConeToneVariable{T}\sim \ConeToneVariable{t_{24}} $};
\node[below,font=\scriptsize] at (0,0) {$0$};\end{tikzpicture}\hfill
\begin{tikzpicture}[x=0.78cm,y=2.8cm]
\fill[COne Tone probability area] (-0.95,0) -- plot[domain=-0.95000000:1.05000000,samples=100] (\x,{0.3989422804*exp(-\x*\x/2)}) -- (1.05,0) -- cycle;
\draw[thick,COne Tone density curve] plot[domain=-4.50000000:4.50000000,samples=100] (\x,{0.3989422804*exp(-\x*\x/2)});
\draw[->] (-4.7,0)--(4.8,0);
\draw[densely dashed] (-0.95,0)--(-0.95,0.2540591);\node[below,font=\scriptsize] at (-0.95,0) {$ -0.950 $};
\draw[densely dashed] (1.05,0)--(1.05,0.2298821);\node[below,font=\scriptsize] at (1.05,0) {$ 1.050 $};
\node[font=\small] at (0,0.49) {$ \ConeToneVariable{Z}\sim N(0,1) $};
\node[below,font=\scriptsize] at (0,0) {$0$};\end{tikzpicture}
\end{center}
For a negative boundary, use $\ConeToneCDF{F_\nu}(-\ConeToneVariable{c})=\ConeToneCDF{1}-\ConeToneCDF{F_\nu}(\ConeToneVariable{c})$ (and $\ConeToneCDF{\Phi}(-\ConeToneVariable{c})=\ConeToneCDF{1}-\ConeToneCDF{\Phi}(\ConeToneVariable{c})$). Subtract the left cumulative area from the right cumulative area.
By symmetry, $\ConeToneCDF{F_{24}}(-0.95)=\ConeToneCDF{1}-\ConeToneCDF{F_{24}}(0.95)\approx \ConeToneCDF{1}-0.8242092=\ConeToneCDF{0.1757908}$. Supplementary CDF values: $\ConeToneCDF{F_{24}}(-0.95)\approx \ConeToneCDF{0.1757908},\quad \ConeToneCDF{F_{24}}(1.05)\approx \ConeToneCDF{0.8479116}$.
\[\ConeToneProbability{P}(-0.95<\ConeToneVariable{T}<1.05)=\ConeToneCDF{F_{24}}(\ConeToneProbability{1.05})-\ConeToneCDF{F_{24}}(-0.95)\approx \ConeToneProbability{0.6721}.\]
\[\ConeToneProbability{P}(-0.95<\ConeToneVariable{Z}<1.05)=\ConeToneCDF{\Phi}(\ConeToneProbability{1.05})-\ConeToneCDF{\Phi}(-0.95)\approx \ConeToneProbability{0.6821}.\]
\[\CsixToneAbsoluteError{A}=|\CsixToneAbsoluteError{0.6820848}-(0.6721208)|\approx \CsixToneAbsoluteError{0.009964},\quad \CsixToneRelativeError{R}=\frac{\CsixToneAbsoluteError{A}}{\CeightToneScale{0.6721208}}\approx \CsixToneRelativeError{0.014825},\quad \CsixToneRelativeError{E}\approx \CsixToneRelativeError{1.48\%}.\]
\textbf{Interpretation.} Under t, the probability that a commute lasts between 9.5 minutes below and 10.5 minutes above forecast is 67.21\%. The normal approximation overestimates that probability, by 1.00 percentage points (1.48\% relative to t).

\secInicio[c9-d5][sec:c9-direct]

\clearpage
\subsection{Revenue planning alerts}
\label{c9-d6}
\textbf{Problem.} A business models weekly revenue deviations to set planning alerts. Negative deviations mean revenue below plan. Let $X$ be weekly revenue minus its plan, measured in thousands of monetary units. For this exercise, assume the location-scale model $X=0+8T$, where $T\sim t_{5}$ has 5 degrees of freedom. The location is $m=0$ and the scale is $a=8$ thousands of monetary units; $a$ is a scale parameter, not automatically a standard deviation. Compare the reference model $X_T=0+8T$ with the normal approximation $X_N=0+8Z$, $Z\sim N(0,1)$, using identical location and scale.
\QuestionsTasks[sec:c9-direct]{{9}/{sol:c9-d6}/{Find the probability that revenue is more than 8,400 monetary units below plan or more than 10,800 above plan. Also find the probability that revenue differs from plan by more than 16,000 monetary units. Standardize the boundaries, calculate the requested probabilities, and interpret each answer. Solve the same task under both models. Compare each probability or each contextual cutoff separately using absolute error and percentage error relative to t, calculated before rounding.}}
\subsubsection*{Worked solution}\phantomsection\label{sol:c9-d6}
\par\textbf{Translating the information into mathematics}\par
Let $\ConeToneVariable{X_T}$ and $\ConeToneVariable{X_N}$ denote weekly revenue minus its plan, in units of 1,000 monetary units. The stated models specify $\ConeToneVariable{T}\sim \ConeToneVariable{t_{5}}$ and $\ConeToneVariable{Z}\sim N(0,1)$, and the location and scale specify $\ConeToneVariable{X_T}=\ConeToneVariable{0}+8\ConeToneVariable{T}$, $\ConeToneVariable{X_N}=\ConeToneVariable{0}+8\ConeToneVariable{Z}$. Thus $\CfiveTonePopulationMean{m}=\CfiveTonePopulationMean{0}$ and $\CeightToneScale{a}=\CeightToneScale{8}$; the scale is not a stipulated standard deviation. A contextual boundary $\ConeToneVariable{x}$ becomes $\ConeToneVariable{c}=(\ConeToneVariable{x}\ConeToneVariable{-0})/8$. The requested events are $\ConeToneProbability{P}(\ConeToneVariable{X}<-8.4\text{ or }\ConeToneVariable{X}>10.8),\ \ConeToneProbability{P}(\ConeToneVariable{X}<-16\text{ or }\ConeToneVariable{X}>16)$, with $\ConeToneVariable{X}=\ConeToneVariable{X_T}$ or $\ConeToneVariable{X_N}$ for the appropriate model. Write the normal result as $\ConeToneProbability{\ConeToneVariable{r_N}}$ and the t reference as $\ConeToneProbability{\ConeToneVariable{r_T}}$. Compute $\CsixToneAbsoluteError{A}=|\ConeToneProbability{\ConeToneVariable{r_N}}-\ConeToneProbability{\ConeToneVariable{r_T}}|$, $\CsixToneRelativeError{R}=\CsixToneAbsoluteError{A}/|\ConeToneProbability{\ConeToneVariable{r_T}}|$, and $\CsixToneRelativeError{E}=\CsixToneRelativeError{100}\CsixToneRelativeError{R}\%$ with unrounded values.
\par\textbf{Mathematical reasoning and calculations}\par
Standardize each boundary using $(\ConeToneVariable{x}-0)/8$. In the task, $\ConeToneVariable{X}$ means $\ConeToneVariable{X_T}$ for the t model and $\ConeToneVariable{X_N}$ for the normal approximation.
\par\textbf{Part 1.}\[\ConeToneProbability{P}(\ConeToneVariable{X_T}<-8.4\ \text{or}\ \ConeToneVariable{X_T}>10.8)=\ConeToneProbability{P}(\ConeToneVariable{T}<\ConeToneProbability{-1.05}\ \text{or}\ \ConeToneVariable{T}>1.35).\]
\begin{center}
\begin{tikzpicture}[x=0.78cm,y=2.8cm]
\fill[COne Tone probability area] (-4.5,0) -- plot[domain=-4.50000000:-1.05000000,samples=100] (\x,{0.3796066898*(1+\x*\x/5)^(-3)}) -- (-1.05,0) -- cycle;
\fill[COne Tone probability area] (1.35,0) -- plot[domain=1.35000000:4.50000000,samples=100] (\x,{0.3796066898*(1+\x*\x/5)^(-3)}) -- (4.5,0) -- cycle;
\draw[thick,COne Tone density curve] plot[domain=-4.50000000:4.50000000,samples=100] (\x,{0.3796066898*(1+\x*\x/5)^(-3)});
\draw[->] (-4.7,0)--(4.8,0);
\draw[densely dashed] (-1.05,0)--(-1.05,0.2087952);\node[below,font=\scriptsize] at (-1.05,0) {$ -1.050 $};
\draw[densely dashed] (1.35,0)--(1.35,0.1494216);\node[below,font=\scriptsize] at (1.35,0) {$ 1.350 $};
\node[font=\small] at (0,0.49) {$ \ConeToneVariable{T}\sim \ConeToneVariable{t_{5}} $};
\node[below,font=\scriptsize] at (0,0) {$0$};\end{tikzpicture}\hfill
\begin{tikzpicture}[x=0.78cm,y=2.8cm]
\fill[COne Tone probability area] (-4.5,0) -- plot[domain=-4.50000000:-1.05000000,samples=100] (\x,{0.3989422804*exp(-\x*\x/2)}) -- (-1.05,0) -- cycle;
\fill[COne Tone probability area] (1.35,0) -- plot[domain=1.35000000:4.50000000,samples=100] (\x,{0.3989422804*exp(-\x*\x/2)}) -- (4.5,0) -- cycle;
\draw[thick,COne Tone density curve] plot[domain=-4.50000000:4.50000000,samples=100] (\x,{0.3989422804*exp(-\x*\x/2)});
\draw[->] (-4.7,0)--(4.8,0);
\draw[densely dashed] (-1.05,0)--(-1.05,0.2298821);\node[below,font=\scriptsize] at (-1.05,0) {$ -1.050 $};
\draw[densely dashed] (1.35,0)--(1.35,0.1603833);\node[below,font=\scriptsize] at (1.35,0) {$ 1.350 $};
\node[font=\small] at (0,0.49) {$ \ConeToneVariable{Z}\sim N(0,1) $};
\node[below,font=\scriptsize] at (0,0) {$0$};\end{tikzpicture}
\end{center}
For a negative boundary, use $\ConeToneCDF{F_\nu}(-\ConeToneVariable{c})=\ConeToneCDF{1}-\ConeToneCDF{F_\nu}(\ConeToneVariable{c})$ (and $\ConeToneCDF{\Phi}(-\ConeToneVariable{c})=\ConeToneCDF{1}-\ConeToneCDF{\Phi}(\ConeToneVariable{c})$). Add the disjoint outside regions, or subtract the intervening interval from 1.
By symmetry, $\ConeToneCDF{F_{5}}(-1.05)=\ConeToneCDF{1}-\ConeToneCDF{F_{5}}(1.05)\approx \ConeToneCDF{1}-0.8291023=\ConeToneCDF{0.1708977}$. Supplementary CDF values: $\ConeToneCDF{F_{5}}(-1.05)\approx \ConeToneCDF{0.1708977},\quad \ConeToneCDF{F_{5}}(1.35)\approx \ConeToneCDF{0.8825456}$.
\[\ConeToneProbability{P}(\ConeToneVariable{T}<-1.05\ \text{or}\ \ConeToneVariable{T}>1.35)=\ConeToneCDF{F_{5}}(\ConeToneProbability{-1.05})+1-\ConeToneCDF{F_{5}}(1.35)\approx \ConeToneProbability{0.2884}.\]
\[\ConeToneProbability{P}(\ConeToneVariable{Z}<-1.05\ \text{or}\ \ConeToneVariable{Z}>1.35)=\ConeToneCDF{\Phi}(\ConeToneProbability{-1.05})+1-\ConeToneCDF{\Phi}(1.35)\approx \ConeToneProbability{0.2354}.\]
\[\CsixToneAbsoluteError{A}=|\CsixToneAbsoluteError{0.2353670}-(0.2883521)|\approx \CsixToneAbsoluteError{0.052985},\quad \CsixToneRelativeError{R}=\frac{\CsixToneAbsoluteError{A}}{\CeightToneScale{0.2883521}}\approx \CsixToneRelativeError{0.183751},\quad \CsixToneRelativeError{E}\approx \CsixToneRelativeError{18.38\%}.\]
\textbf{Interpretation.} Under t, the probability that weekly revenue is more than 8,400 monetary units below or 10,800 above plan is 28.84\%. The normal approximation underestimates that probability, by 5.30 percentage points (18.38\% relative to t).
\par\textbf{Part 2.}\[\ConeToneProbability{P}(\ConeToneVariable{X_T}<-16\ \text{or}\ \ConeToneVariable{X_T}>16)=\ConeToneProbability{P}(\ConeToneVariable{T}<\ConeToneProbability{-2}\ \text{or}\ \ConeToneVariable{T}>2).\]
\begin{center}
\begin{tikzpicture}[x=0.78cm,y=2.8cm]
\fill[COne Tone probability area] (-4.5,0) -- plot[domain=-4.50000000:-2.00000000,samples=100] (\x,{0.3796066898*(1+\x*\x/5)^(-3)}) -- (-2,0) -- cycle;
\fill[COne Tone probability area] (2,0) -- plot[domain=2.00000000:4.50000000,samples=100] (\x,{0.3796066898*(1+\x*\x/5)^(-3)}) -- (4.5,0) -- cycle;
\draw[thick,COne Tone density curve] plot[domain=-4.50000000:4.50000000,samples=100] (\x,{0.3796066898*(1+\x*\x/5)^(-3)});
\draw[->] (-4.7,0)--(4.8,0);
\draw[densely dashed] (-2,0)--(-2,0.0650903);\node[below,font=\scriptsize] at (-2,0) {$ -2.000 $};
\draw[densely dashed] (2,0)--(2,0.0650903);\node[below,font=\scriptsize] at (2,0) {$ 2.000 $};
\node[font=\small] at (0,0.49) {$ \ConeToneVariable{T}\sim \ConeToneVariable{t_{5}} $};
\node[below,font=\scriptsize] at (0,0) {$0$};\end{tikzpicture}\hfill
\begin{tikzpicture}[x=0.78cm,y=2.8cm]
\fill[COne Tone probability area] (-4.5,0) -- plot[domain=-4.50000000:-2.00000000,samples=100] (\x,{0.3989422804*exp(-\x*\x/2)}) -- (-2,0) -- cycle;
\fill[COne Tone probability area] (2,0) -- plot[domain=2.00000000:4.50000000,samples=100] (\x,{0.3989422804*exp(-\x*\x/2)}) -- (4.5,0) -- cycle;
\draw[thick,COne Tone density curve] plot[domain=-4.50000000:4.50000000,samples=100] (\x,{0.3989422804*exp(-\x*\x/2)});
\draw[->] (-4.7,0)--(4.8,0);
\draw[densely dashed] (-2,0)--(-2,0.0539910);\node[below,font=\scriptsize] at (-2,0) {$ -2.000 $};
\draw[densely dashed] (2,0)--(2,0.0539910);\node[below,font=\scriptsize] at (2,0) {$ 2.000 $};
\node[font=\small] at (0,0.49) {$ \ConeToneVariable{Z}\sim N(0,1) $};
\node[below,font=\scriptsize] at (0,0) {$0$};\end{tikzpicture}
\end{center}
For a negative boundary, use $\ConeToneCDF{F_\nu}(-\ConeToneVariable{c})=\ConeToneCDF{1}-\ConeToneCDF{F_\nu}(\ConeToneVariable{c})$ (and $\ConeToneCDF{\Phi}(-\ConeToneVariable{c})=\ConeToneCDF{1}-\ConeToneCDF{\Phi}(\ConeToneVariable{c})$). Add the disjoint outside regions, or subtract the intervening interval from 1.
By symmetry, $\ConeToneCDF{F_{5}}(-2)=\ConeToneCDF{1}-\ConeToneCDF{F_{5}}(2)\approx \ConeToneCDF{1}-0.9490303=\ConeToneCDF{0.0509697}$. Supplementary CDF values: $\ConeToneCDF{F_{5}}(-2)\approx \ConeToneCDF{0.0509697},\quad \ConeToneCDF{F_{5}}(2)\approx \ConeToneCDF{0.9490303}$.
\[\ConeToneProbability{P}(\ConeToneVariable{T}<-2\ \text{or}\ \ConeToneVariable{T}>2)=\ConeToneCDF{F_{5}}(\ConeToneProbability{-2})+1-\ConeToneCDF{F_{5}}(2)\approx \ConeToneProbability{0.1019}.\]
\[\ConeToneProbability{P}(\ConeToneVariable{Z}<-2\ \text{or}\ \ConeToneVariable{Z}>2)=\ConeToneCDF{\Phi}(\ConeToneProbability{-2})+1-\ConeToneCDF{\Phi}(2)\approx \ConeToneProbability{0.0455}.\]
\[\CsixToneAbsoluteError{A}=|\CsixToneAbsoluteError{0.0455003}-(0.1019395)|\approx \CsixToneAbsoluteError{0.056439},\quad \CsixToneRelativeError{R}=\frac{\CsixToneAbsoluteError{A}}{\CeightToneScale{0.1019395}}\approx \CsixToneRelativeError{0.553654},\quad \CsixToneRelativeError{E}\approx \CsixToneRelativeError{55.37\%}.\]
\textbf{Interpretation.} Under t, the probability that weekly revenue differs from plan by more than 16,000 monetary units is 10.19\%. The normal approximation underestimates that probability, by 5.64 percentage points (55.37\% relative to t).

\secInicio[c9-d6][sec:c9-direct]

\clearpage\section{Inverse probability problems}\label{sec:c9-inverse}
\SectionProblemsTableOfContents{
\SectionProblemLink{c9-i1}{Top delivery delays} \\
\SectionProblemLink{c9-i2}{Bottom cost deviations} \\
\SectionProblemLink{c9-i3}{Central energy deviations} \\
\SectionProblemLink{c9-i4}{Central sales deviations} \\
\SectionProblemLink{c9-i5}{2-tail checks on filling weights} \\
\SectionProblemLink{c9-i6}{2-tail checks on processing times}}

\clearpage
\subsection{Top delivery delays}
\label{c9-i1}
\textbf{Problem.} A courier uses a heavy-tailed planning model for delivery schedule deviations. Negative deviations mean early arrivals. Let $X$ be delivery time minus its scheduled time, measured in minutes. For this exercise, assume the location-scale model $X=0+5T$, where $T\sim t_{9}$ has 9 degrees of freedom. The location is $m=0$ and the scale is $a=5$ minutes; $a$ is a scale parameter, not automatically a standard deviation. Compare the reference model $X_T=0+5T$ with the normal approximation $X_N=0+5Z$, $Z\sim N(0,1)$, using identical location and scale.
\QuestionsTasks[sec:c9-inverse]{{9}/{sol:c9-i1}/{Find the lateness cutoff exceeded by only 5\% of deliveries. State the cumulative probability for each cutoff, use the inverse CDF, and back-transform to the original units. Use both models. Compare corresponding contextual cutoffs separately using absolute and percentage errors relative to t before rounding. Interpret the results.}}
\subsubsection*{Worked solution}\phantomsection\label{sol:c9-i1}
\par\textbf{Translating the information into mathematics}
The reference model is $\ConeToneVariable{X_T}=\ConeToneVariable{5}\ConeToneVariable{T}$, $\ConeToneVariable{T}\sim \ConeToneVariable{t_{9}}$, and the approximation is $\ConeToneVariable{X_N}=\ConeToneVariable{5}\ConeToneVariable{Z}$, $\ConeToneVariable{Z}\sim N(0,1)$. In both cases, standardize a deviation $\ConeToneVariable{x}$ as $\ConeToneVariable{x}/5$. Back-transform each quantile $\ConeToneVariable{c}$ using $\ConeToneVariable{x}=\ConeToneVariable{5}\ConeToneVariable{c}$ (minutes).
\par\textbf{Mathematical reasoning and calculations}\par
\[\ConeToneProbability{P}(\ConeToneVariable{X_T}>\ConeToneVariable{x})=\ConeToneProbability{0.0500}\quad\Longrightarrow\quad \ConeToneProbability{P}(\ConeToneVariable{T}\le \ConeToneVariable{c})=\ConeToneProbability{0.9500},\qquad \ConeToneVariable{c}=\frac{\ConeToneVariable{x}\ConeToneVariable{-0}}{\CeightToneScale{5}}.\]
A top percentage is a right tail, so subtract it from 1 before inverse CDF.
\[\ConeToneVariable{c_T}=\ConeToneCDF{F_{9}^{-1}}(\ConeToneVariable{0.95})\approx \ConeToneVariable{1.833},\qquad \ConeToneVariable{c_N}=\ConeToneCDF{\Phi^{-1}}(\ConeToneVariable{0.95})\approx \ConeToneVariable{1.645}.\]
\begin{center}
\begin{tikzpicture}[x=0.78cm,y=2.8cm]
\fill[COne Tone probability area] (1.8331129326562428,0) -- plot[domain=1.83311293:4.50000000,samples=100] (\x,{0.3880349089*(1+\x*\x/9)^(-5)}) -- (4.5,0) -- cycle;
\draw[thick,COne Tone density curve] plot[domain=-4.50000000:4.50000000,samples=100] (\x,{0.3880349089*(1+\x*\x/9)^(-5)});
\draw[->] (-4.7,0)--(4.8,0);
\draw[densely dashed] (1.8331129326562428,0)--(1.8331129326562428,0.0794214);\node[below,font=\scriptsize] at (1.8331129326562428,0) {$ 1.833 $};
\node[font=\small] at (0,0.49) {$ \ConeToneVariable{T}\sim \ConeToneVariable{t_{9}} $};
\node[below,font=\scriptsize] at (0,0) {$0$};\end{tikzpicture}\hfill
\begin{tikzpicture}[x=0.78cm,y=2.8cm]
\fill[COne Tone probability area] (1.6448536269514715,0) -- plot[domain=1.64485363:4.50000000,samples=100] (\x,{0.3989422804*exp(-\x*\x/2)}) -- (4.5,0) -- cycle;
\draw[thick,COne Tone density curve] plot[domain=-4.50000000:4.50000000,samples=100] (\x,{0.3989422804*exp(-\x*\x/2)});
\draw[->] (-4.7,0)--(4.8,0);
\draw[densely dashed] (1.6448536269514715,0)--(1.6448536269514715,0.1031356);\node[below,font=\scriptsize] at (1.6448536269514715,0) {$ 1.645 $};
\node[font=\small] at (0,0.49) {$ \ConeToneVariable{Z}\sim N(0,1) $};
\node[below,font=\scriptsize] at (0,0) {$0$};\end{tikzpicture}
\end{center}
\[\ConeToneVariable{x_T}=\ConeToneVariable{0}+5\ConeToneVariable{c_T}\approx \ConeToneVariable{9.17}\text{ minutes},\qquad \ConeToneVariable{x_N}=\ConeToneVariable{0}+5\ConeToneVariable{c_N}\approx \ConeToneVariable{8.22}\text{ minutes}.\]
\textbf{Contextual cutoff error.} With unrounded $\ConeToneVariable{\ConeToneVariable{r_T}}=5\ConeToneVariable{c_T}$ and $\ConeToneVariable{\ConeToneVariable{r_N}}=5\ConeToneVariable{c_N}$,
\[\CsixToneAbsoluteError{A}=|\ConeToneVariable{\ConeToneVariable{r_N}}-\ConeToneVariable{\ConeToneVariable{r_T}}|\approx \CsixToneAbsoluteError{0.941297}\text{ minutes},\qquad \CsixToneRelativeError{E}=\CsixToneRelativeError{100}\frac{\CsixToneAbsoluteError{A}}{|\ConeToneVariable{\ConeToneVariable{r_T}}|}\approx \CsixToneRelativeError{10.27\%}.\]
\textbf{Interpretation.} About 5\% of deviations are above 9.17 minutes. Normal places the cutoff closer to the location, so it understates the magnitude of this extreme deviation.

\secInicio[c9-i1][sec:c9-inverse]

\clearpage
\subsection{Bottom cost deviations}
\label{c9-i2}
\textbf{Problem.} A manufacturer models batch cost deviations for budget planning. Negative deviations represent savings. Let $X$ be batch cost minus its budget, measured in monetary units. For this exercise, assume the location-scale model $X=0+20T$, where $T\sim t_{11}$ has 11 degrees of freedom. The location is $m=0$ and the scale is $a=20$ monetary units; $a$ is a scale parameter, not automatically a standard deviation. Compare the reference model $X_T=0+20T$ with the normal approximation $X_N=0+20Z$, $Z\sim N(0,1)$, using identical location and scale.
\QuestionsTasks[sec:c9-inverse]{{9}/{sol:c9-i2}/{Find the cost deviation below which 10\% of batches fall. State the cumulative probability for each cutoff, use the inverse CDF, and back-transform to the original units. Use both models. Compare corresponding contextual cutoffs separately using absolute and percentage errors relative to t before rounding. Interpret the results.}}
\subsubsection*{Worked solution}\phantomsection\label{sol:c9-i2}
\par\textbf{Translating the information into mathematics}
The reference model is $\ConeToneVariable{X_T}=\ConeToneVariable{20}\ConeToneVariable{T}$, $\ConeToneVariable{T}\sim \ConeToneVariable{t_{11}}$, and the approximation is $\ConeToneVariable{X_N}=\ConeToneVariable{20}\ConeToneVariable{Z}$, $\ConeToneVariable{Z}\sim N(0,1)$. In both cases, standardize a deviation $\ConeToneVariable{x}$ as $\ConeToneVariable{x}/20$. Back-transform each quantile $\ConeToneVariable{c}$ using $\ConeToneVariable{x}=\ConeToneVariable{20}\ConeToneVariable{c}$ (monetary units).
\par\textbf{Mathematical reasoning and calculations}\par
\[\ConeToneProbability{P}(\ConeToneVariable{X_T}<\ConeToneVariable{x})=\ConeToneProbability{0.1000}\quad\Longrightarrow\quad \ConeToneProbability{P}(\ConeToneVariable{T}\le \ConeToneVariable{c})=\ConeToneProbability{0.1000},\qquad \ConeToneVariable{c}=\frac{\ConeToneVariable{x}\ConeToneVariable{-0}}{\CeightToneScale{20}}.\]
The bottom percentage is already a left cumulative probability; the cutoff is negative by symmetry.
\[\ConeToneVariable{c_T}=\ConeToneCDF{F_{11}^{-1}}(\ConeToneVariable{0.1})\approx \ConeToneVariable{-1.363},\qquad \ConeToneVariable{c_N}=\ConeToneCDF{\Phi^{-1}}(\ConeToneVariable{0.1})\approx \ConeToneVariable{-1.282}.\]
\begin{center}
\begin{tikzpicture}[x=0.78cm,y=2.8cm]
\fill[COne Tone probability area] (-4.5,0) -- plot[domain=-4.50000000:-1.36343032,samples=100] (\x,{0.3899897571*(1+\x*\x/11)^(-6)}) -- (-1.363430318020539,0) -- cycle;
\draw[thick,COne Tone density curve] plot[domain=-4.50000000:4.50000000,samples=100] (\x,{0.3899897571*(1+\x*\x/11)^(-6)});
\draw[->] (-4.7,0)--(4.8,0);
\draw[densely dashed] (-1.363430318020539,0)--(-1.363430318020539,0.1528192);\node[below,font=\scriptsize] at (-1.363430318020539,0) {$ -1.363 $};
\node[font=\small] at (0,0.49) {$ \ConeToneVariable{T}\sim \ConeToneVariable{t_{11}} $};
\node[below,font=\scriptsize] at (0,0) {$0$};\end{tikzpicture}\hfill
\begin{tikzpicture}[x=0.78cm,y=2.8cm]
\fill[COne Tone probability area] (-4.5,0) -- plot[domain=-4.50000000:-1.28155157,samples=100] (\x,{0.3989422804*exp(-\x*\x/2)}) -- (-1.2815515655446008,0) -- cycle;
\draw[thick,COne Tone density curve] plot[domain=-4.50000000:4.50000000,samples=100] (\x,{0.3989422804*exp(-\x*\x/2)});
\draw[->] (-4.7,0)--(4.8,0);
\draw[densely dashed] (-1.2815515655446008,0)--(-1.2815515655446008,0.1754983);\node[below,font=\scriptsize] at (-1.2815515655446008,0) {$ -1.282 $};
\node[font=\small] at (0,0.49) {$ \ConeToneVariable{Z}\sim N(0,1) $};
\node[below,font=\scriptsize] at (0,0) {$0$};\end{tikzpicture}
\end{center}
\[\ConeToneVariable{x_T}=\ConeToneVariable{0}+20\ConeToneVariable{c_T}\approx \ConeToneVariable{-27.27}\text{ monetary units},\qquad \ConeToneVariable{x_N}=\ConeToneVariable{0}+20\ConeToneVariable{c_N}\approx \ConeToneVariable{-25.63}\text{ monetary units}.\]
\textbf{Contextual cutoff error.} With unrounded $\ConeToneVariable{\ConeToneVariable{r_T}}=20\ConeToneVariable{c_T}$ and $\ConeToneVariable{\ConeToneVariable{r_N}}=20\ConeToneVariable{c_N}$,
\[\CsixToneAbsoluteError{A}=|\ConeToneVariable{\ConeToneVariable{r_N}}-\ConeToneVariable{\ConeToneVariable{r_T}}|\approx \CsixToneAbsoluteError{1.637575}\text{ monetary units},\qquad \CsixToneRelativeError{E}=\CsixToneRelativeError{100}\frac{\CsixToneAbsoluteError{A}}{|\ConeToneVariable{\ConeToneVariable{r_T}}|}\approx \CsixToneRelativeError{6.01\%}.\]
\textbf{Interpretation.} About 10\% of deviations are below -27.27 monetary units. Normal places the cutoff closer to the location, so it understates the magnitude of this extreme deviation.

\secInicio[c9-i2][sec:c9-inverse]

\clearpage
\subsection{Central energy deviations}
\label{c9-i3}
\textbf{Problem.} An energy planner models household electricity forecast errors. Negative errors mean use below forecast. Let $X$ be household electricity use minus its forecast, measured in kilowatt-hours. For this exercise, assume the location-scale model $X=0+50T$, where $T\sim t_{19}$ has 19 degrees of freedom. The location is $m=0$ and the scale is $a=50$ kilowatt-hours; $a$ is a scale parameter, not automatically a standard deviation. Compare the reference model $X_T=0+50T$ with the normal approximation $X_N=0+50Z$, $Z\sim N(0,1)$, using identical location and scale.
\QuestionsTasks[sec:c9-inverse]{{9}/{sol:c9-i3}/{Find two electricity-use deviation cutoffs enclosing the middle 95\% of outcomes, with equal probabilities in the two excluded tails. State the cumulative probability for each cutoff, use the inverse CDF, and back-transform to the original units. Use both models. Compare corresponding contextual cutoffs separately using absolute and percentage errors relative to t before rounding. Interpret the results.}}
\subsubsection*{Worked solution}\phantomsection\label{sol:c9-i3}
\par\textbf{Translating the information into mathematics}
The reference model is $\ConeToneVariable{X_T}=\ConeToneVariable{50}\ConeToneVariable{T}$, $\ConeToneVariable{T}\sim \ConeToneVariable{t_{19}}$, and the approximation is $\ConeToneVariable{X_N}=\ConeToneVariable{50}\ConeToneVariable{Z}$, $\ConeToneVariable{Z}\sim N(0,1)$. In both cases, standardize a deviation $\ConeToneVariable{x}$ as $\ConeToneVariable{x}/50$. Back-transform each quantile $\ConeToneVariable{c}$ using $\ConeToneVariable{x}=\ConeToneVariable{50}\ConeToneVariable{c}$ (kWh). Equal tails give symmetric contextual limits $\CeightToneBounds{L}=\CeightToneBounds{-50}\ConeToneVariable{c}$ and $\CeightToneBounds{U}=\CeightToneBounds{50}\ConeToneVariable{c}$.
\par\textbf{Mathematical reasoning and calculations}\par
\[\ConeToneProbability{P}(-\ConeToneVariable{c}<\ConeToneVariable{T}<\ConeToneVariable{c})=\ConeToneProbability{0.9500}\quad\Longrightarrow\quad \ConeToneProbability{P}(\ConeToneVariable{T}>\ConeToneVariable{c})=\ConeToneProbability{P}(\ConeToneVariable{T}<-\ConeToneVariable{c})=\ConeToneProbability{0.0250}.\]
The lower cutoff has cumulative probability $\ConeToneProbability{0.0250}$; the upper cutoff has cumulative probability $\ConeToneProbability{0.9750}$.
Divide the outside complement $1-\ConeToneProbability{p}$ by 2 because symmetry makes the tails equal.
\[\ConeToneCDF{F_{19}}(\ConeToneVariable{c})=\ConeToneCDF{0.9750},\qquad \ConeToneVariable{c_T}=\ConeToneCDF{F_{19}^{-1}}(\ConeToneVariable{0.975})\approx \ConeToneVariable{2.093},\quad \ConeToneVariable{c_N}=\ConeToneCDF{\Phi^{-1}}(\ConeToneVariable{0.975})\approx \ConeToneVariable{1.960}.\]
\begin{center}
\begin{tikzpicture}[x=0.78cm,y=2.8cm]
\fill[COne Tone probability area] (-2.093024054408329,0) -- plot[domain=-2.09302405:2.09302405,samples=100] (\x,{0.3937298073*(1+\x*\x/19)^(-10)}) -- (2.093024054408329,0) -- cycle;
\draw[thick,COne Tone density curve] plot[domain=-4.50000000:4.50000000,samples=100] (\x,{0.3937298073*(1+\x*\x/19)^(-10)});
\draw[->] (-4.7,0)--(4.8,0);
\draw[densely dashed] (-2.093024054408329,0)--(-2.093024054408329,0.0494481);\node[below,font=\scriptsize] at (-2.093024054408329,0) {$ -2.093 $};
\draw[densely dashed] (2.093024054408329,0)--(2.093024054408329,0.0494481);\node[below,font=\scriptsize] at (2.093024054408329,0) {$ 2.093 $};
\node[font=\small] at (0,0.49) {$ \ConeToneVariable{T}\sim \ConeToneVariable{t_{19}} $};
\node[below,font=\scriptsize] at (0,0) {$0$};\end{tikzpicture}\hfill
\begin{tikzpicture}[x=0.78cm,y=2.8cm]
\fill[COne Tone probability area] (-1.9599639845400536,0) -- plot[domain=-1.95996398:1.95996398,samples=100] (\x,{0.3989422804*exp(-\x*\x/2)}) -- (1.9599639845400536,0) -- cycle;
\draw[thick,COne Tone density curve] plot[domain=-4.50000000:4.50000000,samples=100] (\x,{0.3989422804*exp(-\x*\x/2)});
\draw[->] (-4.7,0)--(4.8,0);
\draw[densely dashed] (-1.9599639845400536,0)--(-1.9599639845400536,0.0584451);\node[below,font=\scriptsize] at (-1.9599639845400536,0) {$ -1.960 $};
\draw[densely dashed] (1.9599639845400536,0)--(1.9599639845400536,0.0584451);\node[below,font=\scriptsize] at (1.9599639845400536,0) {$ 1.960 $};
\node[font=\small] at (0,0.49) {$ \ConeToneVariable{Z}\sim N(0,1) $};
\node[below,font=\scriptsize] at (0,0) {$0$};\end{tikzpicture}
\end{center}
The standardized endpoints are $[-2.093,2.093]$ under t and $[-1.960,1.960]$ under normal.
\[\ConeToneVariable{L_T}=\ConeToneVariable{0}-50\ConeToneVariable{c_T}\approx \ConeToneVariable{-104.65},\qquad \ConeToneVariable{U_T}=\ConeToneVariable{0}+50\ConeToneVariable{c_T}\approx \ConeToneVariable{104.65}\quad(\text{kWh}).\]
Normal gives $[-98.00,98.00]$ kWh.
\textbf{Endpoint comparisons.} All cutoffs and absolute errors below are in kilowatt-hours. Calculate $\CsixToneAbsoluteError{A}=|\ConeToneVariable{\ConeToneVariable{r_N}}-\ConeToneVariable{\ConeToneVariable{r_T}}|$ and $\CsixToneRelativeError{E}=\CsixToneRelativeError{100}\CsixToneAbsoluteError{A}/|\ConeToneVariable{\ConeToneVariable{r_T}}|$ before rounding, separately for each endpoint.
\begin{center}\begin{tabular}{lrrrr}\toprule Endpoint & t reference & Normal & $\CsixToneAbsoluteError{A}$ & $\CsixToneRelativeError{E}$ (\%)\\\midrule
Lower&-104.651203&-97.998199&6.653003&6.36\\
Upper&104.651203&97.998199&6.653003&6.36\\
\bottomrule\end{tabular}\end{center}
\textbf{Interpretation.} The t interval contains 95\% of modeled deviations; 2.5\% lie in each excluded tail. The normal limits are less extreme because normal has lighter tails; using them for t would trigger more alerts or cover less than the intended central percentage.

\secInicio[c9-i3][sec:c9-inverse]

\clearpage
\subsection{Central sales deviations}
\label{c9-i4}
\textbf{Problem.} A retailer uses a continuous model for monthly sales deviations when planning inventory. Negative deviations mean sales below target. Let $X$ be monthly sales minus its target, measured in units. For this exercise, assume the location-scale model $X=0+20T$, where $T\sim t_{29}$ has 29 degrees of freedom. The location is $m=0$ and the scale is $a=20$ units; $a$ is a scale parameter, not automatically a standard deviation. Compare the reference model $X_T=0+20T$ with the normal approximation $X_N=0+20Z$, $Z\sim N(0,1)$, using identical location and scale.
\QuestionsTasks[sec:c9-inverse]{{9}/{sol:c9-i4}/{Find two sales-deviation cutoffs enclosing the middle 80\% of outcomes, with equal probabilities in the two excluded tails. State the cumulative probability for each cutoff, use the inverse CDF, and back-transform to the original units. Use both models. Compare corresponding contextual cutoffs separately using absolute and percentage errors relative to t before rounding. Interpret the results.}}
\subsubsection*{Worked solution}\phantomsection\label{sol:c9-i4}
\par\textbf{Translating the information into mathematics}
The reference model is $\ConeToneVariable{X_T}=\ConeToneVariable{20}\ConeToneVariable{T}$, $\ConeToneVariable{T}\sim \ConeToneVariable{t_{29}}$, and the approximation is $\ConeToneVariable{X_N}=\ConeToneVariable{20}\ConeToneVariable{Z}$, $\ConeToneVariable{Z}\sim N(0,1)$. In both cases, standardize a deviation $\ConeToneVariable{x}$ as $\ConeToneVariable{x}/20$. Back-transform each quantile $\ConeToneVariable{c}$ using $\ConeToneVariable{x}=\ConeToneVariable{20}\ConeToneVariable{c}$ (units). Equal tails give symmetric contextual limits $\CeightToneBounds{L}=\CeightToneBounds{-20}\ConeToneVariable{c}$ and $\CeightToneBounds{U}=\CeightToneBounds{20}\ConeToneVariable{c}$.
\par\textbf{Mathematical reasoning and calculations}\par
\[\ConeToneProbability{P}(-\ConeToneVariable{c}<\ConeToneVariable{T}<\ConeToneVariable{c})=\ConeToneProbability{0.8000}\quad\Longrightarrow\quad \ConeToneProbability{P}(\ConeToneVariable{T}>\ConeToneVariable{c})=\ConeToneProbability{P}(\ConeToneVariable{T}<-\ConeToneVariable{c})=\ConeToneProbability{0.1000}.\]
The lower cutoff has cumulative probability $\ConeToneProbability{0.1000}$; the upper cutoff has cumulative probability $\ConeToneProbability{0.9000}$.
Divide the outside complement $1-\ConeToneProbability{p}$ by 2 because symmetry makes the tails equal.
\[\ConeToneCDF{F_{29}}(\ConeToneVariable{c})=\ConeToneCDF{0.9000},\qquad \ConeToneVariable{c_T}=\ConeToneCDF{F_{29}^{-1}}(\ConeToneVariable{0.9})\approx \ConeToneVariable{1.311},\quad \ConeToneVariable{c_N}=\ConeToneCDF{\Phi^{-1}}(\ConeToneVariable{0.9})\approx \ConeToneVariable{1.282}.\]
\begin{center}
\begin{tikzpicture}[x=0.78cm,y=2.8cm]
\fill[COne Tone probability area] (-1.3114336473015578,0) -- plot[domain=-1.31143365:1.31143365,samples=100] (\x,{0.3955185790*(1+\x*\x/29)^(-15)}) -- (1.3114336473015578,0) -- cycle;
\draw[thick,COne Tone density curve] plot[domain=-4.50000000:4.50000000,samples=100] (\x,{0.3955185790*(1+\x*\x/29)^(-15)});
\draw[->] (-4.7,0)--(4.8,0);
\draw[densely dashed] (-1.3114336473015578,0)--(-1.3114336473015578,0.1666667);\node[below,font=\scriptsize] at (-1.3114336473015578,0) {$ -1.311 $};
\draw[densely dashed] (1.3114336473015578,0)--(1.3114336473015578,0.1666667);\node[below,font=\scriptsize] at (1.3114336473015578,0) {$ 1.311 $};
\node[font=\small] at (0,0.49) {$ \ConeToneVariable{T}\sim \ConeToneVariable{t_{29}} $};
\node[below,font=\scriptsize] at (0,0) {$0$};\end{tikzpicture}\hfill
\begin{tikzpicture}[x=0.78cm,y=2.8cm]
\fill[COne Tone probability area] (-1.2815515655446008,0) -- plot[domain=-1.28155157:1.28155157,samples=100] (\x,{0.3989422804*exp(-\x*\x/2)}) -- (1.2815515655446008,0) -- cycle;
\draw[thick,COne Tone density curve] plot[domain=-4.50000000:4.50000000,samples=100] (\x,{0.3989422804*exp(-\x*\x/2)});
\draw[->] (-4.7,0)--(4.8,0);
\draw[densely dashed] (-1.2815515655446008,0)--(-1.2815515655446008,0.1754983);\node[below,font=\scriptsize] at (-1.2815515655446008,0) {$ -1.282 $};
\draw[densely dashed] (1.2815515655446008,0)--(1.2815515655446008,0.1754983);\node[below,font=\scriptsize] at (1.2815515655446008,0) {$ 1.282 $};
\node[font=\small] at (0,0.49) {$ \ConeToneVariable{Z}\sim N(0,1) $};
\node[below,font=\scriptsize] at (0,0) {$0$};\end{tikzpicture}
\end{center}
The standardized endpoints are $[-1.311,1.311]$ under t and $[-1.282,1.282]$ under normal.
\[\ConeToneVariable{L_T}=\ConeToneVariable{0}-20\ConeToneVariable{c_T}\approx \ConeToneVariable{-26.23},\qquad \ConeToneVariable{U_T}=\ConeToneVariable{0}+20\ConeToneVariable{c_T}\approx \ConeToneVariable{26.23}\quad(\text{units}).\]
Normal gives $[-25.63,25.63]$ units.
\textbf{Endpoint comparisons.} All cutoffs and absolute errors below are in units. Calculate $\CsixToneAbsoluteError{A}=|\ConeToneVariable{\ConeToneVariable{r_N}}-\ConeToneVariable{\ConeToneVariable{r_T}}|$ and $\CsixToneRelativeError{E}=\CsixToneRelativeError{100}\CsixToneAbsoluteError{A}/|\ConeToneVariable{\ConeToneVariable{r_T}}|$ before rounding, separately for each endpoint.
\begin{center}\begin{tabular}{lrrrr}\toprule Endpoint & t reference & Normal & $\CsixToneAbsoluteError{A}$ & $\CsixToneRelativeError{E}$ (\%)\\\midrule
Lower&-26.228673&-25.631031&0.597642&2.28\\
Upper&26.228673&25.631031&0.597642&2.28\\
\bottomrule\end{tabular}\end{center}
\textbf{Interpretation.} The t interval contains 80\% of modeled deviations; 10\% lie in each excluded tail. The normal limits are less extreme because normal has lighter tails; using them for t would trigger more alerts or cover less than the intended central percentage.

\secInicio[c9-i4][sec:c9-inverse]

\clearpage
\subsection{2-tail checks on filling weights}
\label{c9-i5}
\textbf{Problem.} A packaging engineer uses a heavy-tailed model for signed filling-weight deviations to set inspection limits. Let $X$ be filling weight minus its target, measured in grams. For this exercise, assume the location-scale model $X=0+2T$, where $T\sim t_{5}$ has 5 degrees of freedom. The location is $m=0$ and the scale is $a=2$ grams; $a$ is a scale parameter, not automatically a standard deviation. Compare the reference model $X_T=0+2T$ with the normal approximation $X_N=0+2Z$, $Z\sim N(0,1)$, using identical location and scale.
\QuestionsTasks[sec:c9-inverse]{{9}/{sol:c9-i5}/{Find lower and upper filling-weight deviation cutoffs leaving 10\% of outcomes outside, with 5\% in each tail. State the cumulative probability for each cutoff, use the inverse CDF, and back-transform to the original units. Use both models. Compare corresponding contextual cutoffs separately using absolute and percentage errors relative to t before rounding. Interpret the results.}}
\subsubsection*{Worked solution}\phantomsection\label{sol:c9-i5}
\par\textbf{Translating the information into mathematics}
The reference model is $\ConeToneVariable{X_T}=\ConeToneVariable{2}\ConeToneVariable{T}$, $\ConeToneVariable{T}\sim \ConeToneVariable{t_{5}}$, and the approximation is $\ConeToneVariable{X_N}=\ConeToneVariable{2}\ConeToneVariable{Z}$, $\ConeToneVariable{Z}\sim N(0,1)$. In both cases, standardize a deviation $\ConeToneVariable{x}$ as $\ConeToneVariable{x}/2$. Back-transform each quantile $\ConeToneVariable{c}$ using $\ConeToneVariable{x}=\ConeToneVariable{2}\ConeToneVariable{c}$ (grams). Equal tails give symmetric contextual limits $\CeightToneBounds{L}=\CeightToneBounds{-2}\ConeToneVariable{c}$ and $\CeightToneBounds{U}=\CeightToneBounds{2}\ConeToneVariable{c}$.
\par\textbf{Mathematical reasoning and calculations}\par
\[\ConeToneProbability{P}(|\ConeToneVariable{T}|>\ConeToneVariable{c})=\ConeToneProbability{0.1000}\quad\Longrightarrow\quad \ConeToneProbability{P}(\ConeToneVariable{T}>\ConeToneVariable{c})=\ConeToneProbability{P}(\ConeToneVariable{T}<-\ConeToneVariable{c})=\ConeToneProbability{0.0500}.\]
The lower cutoff has cumulative probability $\ConeToneProbability{0.0500}$; the upper cutoff has cumulative probability $\ConeToneProbability{0.9500}$.
Divide the combined outside area by 2 because symmetry makes the tails equal.
\[\ConeToneCDF{F_{5}}(\ConeToneVariable{c})=\ConeToneCDF{0.9500},\qquad \ConeToneVariable{c_T}=\ConeToneCDF{F_{5}^{-1}}(\ConeToneVariable{0.95})\approx \ConeToneVariable{2.015},\quad \ConeToneVariable{c_N}=\ConeToneCDF{\Phi^{-1}}(\ConeToneVariable{0.95})\approx \ConeToneVariable{1.645}.\]
\begin{center}
\begin{tikzpicture}[x=0.78cm,y=2.8cm]
\fill[COne Tone probability area] (-4.5,0) -- plot[domain=-4.50000000:-2.01504837,samples=100] (\x,{0.3796066898*(1+\x*\x/5)^(-3)}) -- (-2.0150483733330287,0) -- cycle;
\fill[COne Tone probability area] (2.0150483733330287,0) -- plot[domain=2.01504837:4.50000000,samples=100] (\x,{0.3796066898*(1+\x*\x/5)^(-3)}) -- (4.5,0) -- cycle;
\draw[thick,COne Tone density curve] plot[domain=-4.50000000:4.50000000,samples=100] (\x,{0.3796066898*(1+\x*\x/5)^(-3)});
\draw[->] (-4.7,0)--(4.8,0);
\draw[densely dashed] (-2.0150483733330287,0)--(-2.0150483733330287,0.0637968);\node[below,font=\scriptsize] at (-2.0150483733330287,0) {$ -2.015 $};
\draw[densely dashed] (2.0150483733330287,0)--(2.0150483733330287,0.0637968);\node[below,font=\scriptsize] at (2.0150483733330287,0) {$ 2.015 $};
\node[font=\small] at (0,0.49) {$ \ConeToneVariable{T}\sim \ConeToneVariable{t_{5}} $};
\node[below,font=\scriptsize] at (0,0) {$0$};\end{tikzpicture}\hfill
\begin{tikzpicture}[x=0.78cm,y=2.8cm]
\fill[COne Tone probability area] (-4.5,0) -- plot[domain=-4.50000000:-1.64485363,samples=100] (\x,{0.3989422804*exp(-\x*\x/2)}) -- (-1.6448536269514715,0) -- cycle;
\fill[COne Tone probability area] (1.6448536269514715,0) -- plot[domain=1.64485363:4.50000000,samples=100] (\x,{0.3989422804*exp(-\x*\x/2)}) -- (4.5,0) -- cycle;
\draw[thick,COne Tone density curve] plot[domain=-4.50000000:4.50000000,samples=100] (\x,{0.3989422804*exp(-\x*\x/2)});
\draw[->] (-4.7,0)--(4.8,0);
\draw[densely dashed] (-1.6448536269514715,0)--(-1.6448536269514715,0.1031356);\node[below,font=\scriptsize] at (-1.6448536269514715,0) {$ -1.645 $};
\draw[densely dashed] (1.6448536269514715,0)--(1.6448536269514715,0.1031356);\node[below,font=\scriptsize] at (1.6448536269514715,0) {$ 1.645 $};
\node[font=\small] at (0,0.49) {$ \ConeToneVariable{Z}\sim N(0,1) $};
\node[below,font=\scriptsize] at (0,0) {$0$};\end{tikzpicture}
\end{center}
\[\begin{aligned}\ConeToneVariable{L_T}=\ConeToneVariable{-2}\ConeToneVariable{c_T}\approx \ConeToneVariable{-4.03},\qquad \ConeToneVariable{U_T}=\ConeToneVariable{2}\ConeToneVariable{c_T}\approx \ConeToneVariable{4.03}\quad(\text{grams}).\\L_N=-2\ConeToneVariable{c_N}\approx \ConeToneVariable{-3.29},\qquad \ConeToneVariable{U_N}=\ConeToneVariable{2}\ConeToneVariable{c_N}\approx \ConeToneVariable{3.29}\quad(\text{grams}).\end{aligned}\]
\textbf{Endpoint comparisons.} All cutoffs and absolute errors below are in grams. Calculate $\CsixToneAbsoluteError{A}=|\ConeToneVariable{\ConeToneVariable{r_N}}-\ConeToneVariable{\ConeToneVariable{r_T}}|$ and $\CsixToneRelativeError{E}=\CsixToneRelativeError{100}\CsixToneAbsoluteError{A}/|\ConeToneVariable{\ConeToneVariable{r_T}}|$ before rounding, separately for each endpoint.
\begin{center}\begin{tabular}{lrrrr}\toprule Endpoint & t reference & Normal & $\CsixToneAbsoluteError{A}$ & $\CsixToneRelativeError{E}$ (\%)\\\midrule
Lower&-4.030097&-3.289707&0.740389&18.37\\
Upper&4.030097&3.289707&0.740389&18.37\\
\bottomrule\end{tabular}\end{center}
\textbf{Interpretation.} The reference limits are 4.03 grams below and above target. The t interval contains 90\% of modeled deviations; 5\% lie in each excluded tail. Normal sets narrower limits and would trigger more alerts under the reference t model.

\secInicio[c9-i5][sec:c9-inverse]

\clearpage
\subsection{2-tail checks on processing times}
\label{c9-i6}
\textbf{Problem.} A production planner uses a heavy-tailed model for signed processing-time deviations to set alert limits. Let $X$ be processing time minus its target, measured in minutes. For this exercise, assume the location-scale model $X=0+3T$, where $T\sim t_{24}$ has 24 degrees of freedom. The location is $m=0$ and the scale is $a=3$ minutes; $a$ is a scale parameter, not automatically a standard deviation. Compare the reference model $X_T=0+3T$ with the normal approximation $X_N=0+3Z$, $Z\sim N(0,1)$, using identical location and scale.
\QuestionsTasks[sec:c9-inverse]{{9}/{sol:c9-i6}/{Find lower and upper processing-time deviation cutoffs leaving 4\% of outcomes outside, with 2\% in each tail. State the cumulative probability for each cutoff, use the inverse CDF, and back-transform to the original units. Use both models. Compare corresponding contextual cutoffs separately using absolute and percentage errors relative to t before rounding. Interpret the results.}}
\subsubsection*{Worked solution}\phantomsection\label{sol:c9-i6}
\par\textbf{Translating the information into mathematics}
The reference model is $\ConeToneVariable{X_T}=\ConeToneVariable{3}\ConeToneVariable{T}$, $\ConeToneVariable{T}\sim \ConeToneVariable{t_{24}}$, and the approximation is $\ConeToneVariable{X_N}=\ConeToneVariable{3}\ConeToneVariable{Z}$, $\ConeToneVariable{Z}\sim N(0,1)$. In both cases, standardize a deviation $\ConeToneVariable{x}$ as $\ConeToneVariable{x}/3$. Back-transform each quantile $\ConeToneVariable{c}$ using $\ConeToneVariable{x}=\ConeToneVariable{3}\ConeToneVariable{c}$ (minutes). Equal tails give symmetric contextual limits $\CeightToneBounds{L}=\CeightToneBounds{-3}\ConeToneVariable{c}$ and $\CeightToneBounds{U}=\CeightToneBounds{3}\ConeToneVariable{c}$.
\par\textbf{Mathematical reasoning and calculations}\par
\[\ConeToneProbability{P}(|\ConeToneVariable{T}|>\ConeToneVariable{c})=\ConeToneProbability{0.0400}\quad\Longrightarrow\quad \ConeToneProbability{P}(\ConeToneVariable{T}>\ConeToneVariable{c})=\ConeToneProbability{P}(\ConeToneVariable{T}<-\ConeToneVariable{c})=\ConeToneProbability{0.0200}.\]
The lower cutoff has cumulative probability $\ConeToneProbability{0.0200}$; the upper cutoff has cumulative probability $\ConeToneProbability{0.9800}$.
Divide the combined outside area by 2 because symmetry makes the tails equal.
\[\ConeToneCDF{F_{24}}(\ConeToneVariable{c})=\ConeToneCDF{0.9800},\qquad \ConeToneVariable{c_T}=\ConeToneCDF{F_{24}^{-1}}(\ConeToneVariable{0.98})\approx \ConeToneVariable{2.172},\quad \ConeToneVariable{c_N}=\ConeToneCDF{\Phi^{-1}}(\ConeToneVariable{0.98})\approx \ConeToneVariable{2.054}.\]
\begin{center}
\begin{tikzpicture}[x=0.78cm,y=2.8cm]
\fill[COne Tone probability area] (-4.5,0) -- plot[domain=-4.50000000:-2.17154468,samples=100] (\x,{0.3948093882*(1+\x*\x/24)^(-12.5)}) -- (-2.171544676008055,0) -- cycle;
\fill[COne Tone probability area] (2.171544676008055,0) -- plot[domain=2.17154468:4.50000000,samples=100] (\x,{0.3948093882*(1+\x*\x/24)^(-12.5)}) -- (4.5,0) -- cycle;
\draw[thick,COne Tone density curve] plot[domain=-4.50000000:4.50000000,samples=100] (\x,{0.3948093882*(1+\x*\x/24)^(-12.5)});
\draw[->] (-4.7,0)--(4.8,0);
\draw[densely dashed] (-2.171544676008055,0)--(-2.171544676008055,0.0419327);\node[below,font=\scriptsize] at (-2.171544676008055,0) {$ -2.172 $};
\draw[densely dashed] (2.171544676008055,0)--(2.171544676008055,0.0419327);\node[below,font=\scriptsize] at (2.171544676008055,0) {$ 2.172 $};
\node[font=\small] at (0,0.49) {$ \ConeToneVariable{T}\sim \ConeToneVariable{t_{24}} $};
\node[below,font=\scriptsize] at (0,0) {$0$};\end{tikzpicture}\hfill
\begin{tikzpicture}[x=0.78cm,y=2.8cm]
\fill[COne Tone probability area] (-4.5,0) -- plot[domain=-4.50000000:-2.05374891,samples=100] (\x,{0.3989422804*exp(-\x*\x/2)}) -- (-2.053748910631822,0) -- cycle;
\fill[COne Tone probability area] (2.053748910631822,0) -- plot[domain=2.05374891:4.50000000,samples=100] (\x,{0.3989422804*exp(-\x*\x/2)}) -- (4.5,0) -- cycle;
\draw[thick,COne Tone density curve] plot[domain=-4.50000000:4.50000000,samples=100] (\x,{0.3989422804*exp(-\x*\x/2)});
\draw[->] (-4.7,0)--(4.8,0);
\draw[densely dashed] (-2.053748910631822,0)--(-2.053748910631822,0.0484181);\node[below,font=\scriptsize] at (-2.053748910631822,0) {$ -2.054 $};
\draw[densely dashed] (2.053748910631822,0)--(2.053748910631822,0.0484181);\node[below,font=\scriptsize] at (2.053748910631822,0) {$ 2.054 $};
\node[font=\small] at (0,0.49) {$ \ConeToneVariable{Z}\sim N(0,1) $};
\node[below,font=\scriptsize] at (0,0) {$0$};\end{tikzpicture}
\end{center}
\[\begin{aligned}\ConeToneVariable{L_T}=\ConeToneVariable{-3}\ConeToneVariable{c_T}\approx \ConeToneVariable{-6.51},\qquad \ConeToneVariable{U_T}=\ConeToneVariable{3}\ConeToneVariable{c_T}\approx \ConeToneVariable{6.51}\quad(\text{minutes}).\\L_N=-3\ConeToneVariable{c_N}\approx \ConeToneVariable{-6.16},\qquad \ConeToneVariable{U_N}=\ConeToneVariable{3}\ConeToneVariable{c_N}\approx \ConeToneVariable{6.16}\quad(\text{minutes}).\end{aligned}\]
\textbf{Endpoint comparisons.} All cutoffs and absolute errors below are in minutes. Calculate $\CsixToneAbsoluteError{A}=|\ConeToneVariable{\ConeToneVariable{r_N}}-\ConeToneVariable{\ConeToneVariable{r_T}}|$ and $\CsixToneRelativeError{E}=\CsixToneRelativeError{100}\CsixToneAbsoluteError{A}/|\ConeToneVariable{\ConeToneVariable{r_T}}|$ before rounding, separately for each endpoint.
\begin{center}\begin{tabular}{lrrrr}\toprule Endpoint & t reference & Normal & $\CsixToneAbsoluteError{A}$ & $\CsixToneRelativeError{E}$ (\%)\\\midrule
Lower&-6.514634&-6.161247&0.353387&5.42\\
Upper&6.514634&6.161247&0.353387&5.42\\
\bottomrule\end{tabular}\end{center}
\textbf{Interpretation.} The reference limits are 6.51 minutes below and above target. The t interval contains 96\% of modeled deviations; 2\% lie in each excluded tail. Normal sets narrower limits and would trigger more alerts under the reference t model.

\secInicio[c9-i6][sec:c9-inverse]

\clearpage\section{How sample size improves the approximation}\label{sec:convergence}
\textbf{Problem.} A filling-weight audit takes independent random samples of $\ConeToneUnknown{n}=\ConeToneUnknown{6},20,60,120$ weights from a normal population with mean $\CfiveTonePopulationMean{\mu}$ grams and unknown standard deviation $\CfiveTonePopulationSD{\sigma}$ grams. For each sample, calculate its mean $\CsevenToneSampleMean{\bar X}$ and \CsevenToneSampleSD{sample} \CsevenToneSampleSD{standard} \CsevenToneSampleSD{deviation} $\ConeToneVariable{S}$, both in grams. The audit score is $\ConeToneVariable{T}=(\CsevenToneSampleMean{\bar X}-\CfiveTonePopulationMean{\mu})/(\ConeToneVariable{S}/\sqrt \ConeToneUnknown{n})$, a dimensionless measure of the \CsevenToneSampleMean{sample} \CsevenToneSampleMean{mean}'s deviation from the \CfiveTonePopulationMean{population} \CfiveTonePopulationMean{mean}. Under these assumptions $\ConeToneVariable{T}\sim \ConeToneVariable{t_{n-1}}$. Approximate this score by $\ConeToneVariable{Z}\sim N(0,1)$, with the same location 0 and scale 1. Both $\CsevenToneSampleMean{\bar X}$ and $\ConeToneVariable{S}$ vary between samples.
\QuestionsTasks[sec:convergence]{{9}/{sol:convergence}/{For each sample size, calculate the probability that the audit score exceeds 2 under both models. Then find the score cutoff exceeded by 2.5\% under each model. Compare the probabilities and score cutoffs separately using absolute error and percentage error relative to t, before rounding. Interpret the changes as sample size increases.}}
\subsubsection*{Worked solution}\phantomsection\label{sol:convergence}
\par\textbf{Translating the information into mathematics}
For each \CsixToneCount{sample} \CsixToneCount{size} use $\CeightToneControl{\nu}=\ConeToneUnknown{n}\CeightToneControl{-1}$ and $\ConeToneVariable{T}=(\CsevenToneSampleMean{\bar X}-\CfiveTonePopulationMean{\mu})/(\ConeToneVariable{S}/\sqrt \ConeToneUnknown{n})\sim \ConeToneVariable{t_\nu}$. The normal approximation $\ConeToneVariable{Z}\sim N(0,1)$ has the same location and scale. Compare the same dimensionless audit-score event and target tail area under both models, with t as the reference.
\par\textbf{Mathematical reasoning and calculations}\par
Keep the event fixed: $\ConeToneProbability{P}(\ConeToneVariable{T}>2)=\ConeToneProbability{1}-\ConeToneCDF{F_\nu}(2)$ and $\ConeToneProbability{P}(\ConeToneVariable{Z}>2)=\ConeToneProbability{1}-\ConeToneCDF{\Phi}(2)$.
Keep the inverse target fixed: $\ConeToneProbability{P}(\ConeToneVariable{T}>\ConeToneVariable{c_T})=\ConeToneProbability{P}(\ConeToneVariable{Z}>\ConeToneVariable{c_N})=\ConeToneProbability{0.025}$, so use cumulative area $0.975$.
Here $\CeightToneControl{\nu}=\ConeToneUnknown{n}\CeightToneControl{-1}$; the rows 59 and 119 require calculator supplements.
\begin{center}
\begin{tikzpicture}[x=0.78cm,y=2.8cm]
\fill[COne Tone probability area] (2,0) -- plot[domain=2.00000000:4.50000000,samples=100] (\x,{0.3796066898*(1+\x*\x/5)^(-3)}) -- (4.5,0) -- cycle;
\draw[thick,COne Tone density curve] plot[domain=-4.50000000:4.50000000,samples=100] (\x,{0.3796066898*(1+\x*\x/5)^(-3)});
\draw[->] (-4.7,0)--(4.8,0);
\draw[densely dashed] (2,0)--(2,0.0650903);\node[below,font=\scriptsize] at (2,0) {$ 2.000 $};
\node[font=\small] at (0,0.49) {$ \ConeToneVariable{T}\sim \ConeToneVariable{t_{5}} $};
\node[below,font=\scriptsize] at (0,0) {$0$};\end{tikzpicture}\hfill
\begin{tikzpicture}[x=0.78cm,y=2.8cm]
\fill[COne Tone probability area] (2,0) -- plot[domain=2.00000000:4.50000000,samples=100] (\x,{0.3989422804*exp(-\x*\x/2)}) -- (4.5,0) -- cycle;
\draw[thick,COne Tone density curve] plot[domain=-4.50000000:4.50000000,samples=100] (\x,{0.3989422804*exp(-\x*\x/2)});
\draw[->] (-4.7,0)--(4.8,0);
\draw[densely dashed] (2,0)--(2,0.0539910);\node[below,font=\scriptsize] at (2,0) {$ 2.000 $};
\node[font=\small] at (0,0.49) {$ \ConeToneVariable{Z}\sim N(0,1) $};
\node[below,font=\scriptsize] at (0,0) {$0$};\end{tikzpicture}
\end{center}
\textbf{Probability comparison.} Normal gives $\ConeToneProbability{1}-\ConeToneCDF{\Phi}(\ConeToneProbability{2})\approx \ConeToneCDF{0.0228}$.
\begin{center}\begin{tabular}{rr|rrrrr}\toprule $\ConeToneUnknown{n}$ & $\CeightToneControl{\nu}$ & t result & normal result & $\CsixToneAbsoluteError{A}$ & $\CsixToneRelativeError{R}$ & $\CsixToneRelativeError{E}$ (\%)\\\midrule
6&5&0.0510&0.0228&0.028220&0.553654&55.37\\
20&19&0.0300&0.0228&0.007251&0.241688&24.17\\
60&59&0.0251&0.0228&0.002305&0.092000&9.20\\
120&119&0.0239&0.0228&0.001139&0.047665&4.77\\
\bottomrule\end{tabular}\end{center}
\textbf{Quantile comparison.} Normal gives $\ConeToneCDF{\Phi^{-1}}(0.975)\approx \ConeToneCDF{1.960}$.
\begin{center}\begin{tabular}{rr|rrrrr}\toprule $\ConeToneUnknown{n}$ & $\CeightToneControl{\nu}$ & t result & normal result & $\CsixToneAbsoluteError{A}$ & $\CsixToneRelativeError{R}$ & $\CsixToneRelativeError{E}$ (\%)\\\midrule
6&5&2.571&1.960&0.610618&0.237541&23.75\\
20&19&2.093&1.960&0.133060&0.063573&6.36\\
60&59&2.001&1.960&0.041031&0.020505&2.05\\
120&119&1.980&1.960&0.020136&0.010169&1.02\\
\bottomrule\end{tabular}\end{center}
Probability errors $\CsixToneAbsoluteError{A}$ are on the 0--1 scale; multiply by 100 for percentage-point differences. Cutoff errors are in dimensionless audit-score units. For example, in either table $\CsixToneAbsoluteError{A}=|\ConeToneVariable{\ConeToneVariable{r_N}}-\ConeToneVariable{\ConeToneVariable{r_T}}|$, $\CsixToneRelativeError{R}=\CsixToneAbsoluteError{A}/|\ConeToneVariable{\ConeToneVariable{r_T}}|$, and $\CsixToneRelativeError{E}=\CsixToneRelativeError{100}\CsixToneRelativeError{R}\%$, using unrounded calculator results.
\textbf{Interpretation.} All 3 errors decrease across these sample sizes. Normal underestimates the chance of an audit statistic exceeding 2 and understates the cutoff required for a 2.5\% alert rate. Heavier t tails explain both effects; increasing df reduces the difference. Estimating the \CfiveTonePopulationSD{population} \CfiveTonePopulationSD{standard} \CfiveTonePopulationSD{deviation} still calls for t, even when the approximation is close.
\secInicio[sec:convergence][sec:tools]

\end{document}
